% ============================================================
% THE DIALECTICAL CAGE AND THE GLASS-BOX GOVERNOR
% ============================================================

\documentclass[11pt]{article}

\usepackage{fontspec}
\defaultfontfeatures{Ligatures=TeX,Scale=MatchLowercase}
\setmainfont{Noto Serif}
\setsansfont{Noto Sans}
\setmonofont{DejaVu Sans Mono}
\usepackage{iftex}
\usepackage{microtype}
\ifLuaTeX\microtypesetup{protrusion=true,expansion=true}\else\microtypesetup{protrusion=true,expansion=false}\fi
\usepackage{setspace}
\setstretch{1.06}
\setlength{\parindent}{0pt}
\setlength{\parskip}{0.58em}
\setlength{\emergencystretch}{4.5em}
\clubpenalty=10000
\widowpenalty=10000

\usepackage[
  a4paper,
  left=1.00in,
  right=1.00in,
  top=0.88in,
  bottom=0.92in,
  headheight=22pt,
  headsep=0.26in,
  footskip=0.34in
]{geometry}

\usepackage{amsmath,amssymb,amsthm,mathtools}
\usepackage{unicode-math}
\setmathfont{Latin Modern Math}

\usepackage{enumitem}
\usepackage{array,booktabs,longtable,tabularx}
\usepackage{float}
\usepackage{etoolbox}
\usepackage{needspace}
\setlist{
  leftmargin=*,
  itemsep=0.22em,
  topsep=0.38em,
  parsep=0pt,
  partopsep=0pt
}
\setlist[enumerate]{labelsep=0.55em}
\setlist[itemize]{labelsep=0.55em}

\usepackage{xcolor}
\definecolor{seriesInk}{HTML}{1F2933}
\definecolor{seriesAccent}{HTML}{8A3B4A}
\definecolor{seriesSteel}{HTML}{40566F}
\definecolor{seriesMuted}{HTML}{66727E}
\definecolor{seriesPaper}{HTML}{F7F4EE}
\definecolor{seriesRule}{HTML}{C9D0D6}
\definecolor{stD}{HTML}{2E7A4F}
\definecolor{stN}{HTML}{2D5FA3}
\definecolor{stC}{HTML}{B7791F}
\definecolor{stU}{HTML}{B02A37}

\usepackage{tcolorbox}
\tcbuselibrary{skins,breakable}
\usepackage{tikz}

\usepackage{titlesec}
\titleformat{\section}
  {\normalfont\Large\bfseries\sffamily\color{seriesInk}}
  {\color{seriesAccent}\thesection}
  {0.72em}
  {}
  [\vspace{3pt}{\color{seriesAccent}\titlerule[0.85pt]}]
\titlespacing*{\section}{0pt}{3.0ex plus .8ex minus .2ex}{1.55ex plus .2ex}

\titleformat{name=\section,numberless}
  {\normalfont\Large\bfseries\sffamily\color{seriesInk}}
  {}
  {0pt}
  {}
  [\vspace{3pt}{\color{seriesAccent}\titlerule[0.85pt]}]
\titlespacing*{name=\section,numberless}{0pt}{3.0ex plus .8ex minus .2ex}{1.55ex plus .2ex}

\titleformat{\subsection}
  {\normalfont\large\bfseries\sffamily\color{seriesSteel}}
  {\color{seriesAccent}\thesubsection}
  {0.68em}
  {}
\titlespacing*{\subsection}{0pt}{2.35ex plus .6ex minus .2ex}{0.82ex plus .1ex}

\titleformat{\subsubsection}
  {\normalfont\normalsize\bfseries\sffamily\color{seriesSteel}}
  {\thesubsubsection}
  {0.62em}
  {}
\titlespacing*{\subsubsection}{0pt}{1.8ex plus .4ex minus .2ex}{0.55ex plus .1ex}

\titleformat{name=\subsection,numberless}
  {\normalfont\large\bfseries\sffamily\color{seriesSteel}}
  {}
  {0pt}
  {}
\titlespacing*{name=\subsection,numberless}{0pt}{2.35ex plus .6ex minus .2ex}{0.82ex plus .1ex}

\renewcommand{\sectionmark}[1]{\markboth{\thesection\enspace #1}{}}
\renewcommand{\subsectionmark}[1]{}

\usepackage{fancyhdr}
\usepackage{hyperref}
\usepackage{cleveref}
\usepackage{bookmark}
\usepackage{orcidlink}

\hypersetup{
  colorlinks=true,
  linkcolor=seriesSteel,
  citecolor=seriesAccent,
  urlcolor=seriesAccent,
  filecolor=seriesSteel,
  pdftitle={The Dialectical Cage and the Glass-Box Governor},
  pdfauthor={Canon},
  pdfsubject={A practice-based ethics, safety constitution, and reference-monitor architecture for AI agents},
  pdfkeywords={AI safety, runtime verification, capability security, constitutional governance, neuro-symbolic systems, policy compilation, semantic uncertainty, agent security},
  pdfcreator={LuaLaTeX}
}
\urlstyle{same}

\pdfstringdefDisableCommands{%
  \def\ensuremath#1{#1}%
  \def\mathrm#1{#1}%
  \def\mathsf#1{#1}%
  \def\mathbf#1{#1}%
  \def\mathbb#1{#1}%
  \def\mathcal#1{#1}%
  \def\textsuperscript#1{#1}%
}

\crefname{section}{Section}{Sections}

\newcommand{\statusD}{\textcolor{stD}{\textbf{D}}}
\newcommand{\statusN}{\textcolor{stN}{\textbf{N}}}
\newcommand{\statusC}{\textcolor{stC}{\textbf{C}}}
\newcommand{\statusU}{\textcolor{stU}{\textbf{U}}}
\newcommand{\labelUndetermined}{\textcolor{stU}{\textbf{UNDETERMINED}}}
\newcommand{\labelOutOfScope}{\textcolor{seriesMuted}{\textbf{OUT OF SCOPE}}}
\newcommand{\Atwo}{\texorpdfstring{\ensuremath{\mathrm{A2}'}}{A2'}}
\newcommand{\Ptwo}{\texorpdfstring{\ensuremath{\mathrm{P2}^{*}}}{P2*}}
\newcommand{\Eprime}{\ensuremath{\mathrm{E}'}}
\newcommand{\ADM}{\texorpdfstring{\ensuremath{\mathrm{ADM}_0}}{ADM0}}
\newcommand{\ca}{\texorpdfstring{Class~$\alpha$}{Class alpha}}
\newcommand{\cb}{\texorpdfstring{Class~$\beta$}{Class beta}}
\newcommand{\Gmech}{G-O\textsuperscript{mech}}
\newcommand{\Gsem}{G-O\textsuperscript{sem}}
\newcommand{\CR}{\ensuremath{\mathrm{CR}}}
\newcommand{\CAR}{\ensuremath{\mathrm{CAR}}}
\newcommand{\epscover}{\ensuremath{\epsilon_{\mathrm{cover}}}}
\newcommand{\latTarget}{\ensuremath{L_{p99}}}
\newcommand{\Ground}{\ensuremath{\mathsf{Ground}}}
\newcommand{\DSC}{\ensuremath{\mathsf{DSC}}}
\newcommand{\SafetyBundle}{\textsf{SafetyBundle}}
\newcommand{\Safe}{\ensuremath{\mathsf{Safe}}}
\newcommand{\ALLOW}{\ensuremath{\mathsf{ALLOW}}}
\newcommand{\BLOCK}{\ensuremath{\mathsf{BLOCK}}}
\newcommand{\UNKNOWN}{\ensuremath{\mathsf{UNKNOWN}}}
\newcommand{\PRef}{\ensuremath{\sqsubseteq}}
\newcommand{\ModelHash}{\ensuremath{h_{\mathrm{model}}}}
\newcommand{\PolicyHash}{\ensuremath{h_{\mathrm{policy}}}}
\newcommand{\KernelHash}{\ensuremath{h_{\mathrm{kernel}}}}
\newcommand{\SensorHash}{\ensuremath{h_{\mathrm{sensor}}}}
\newcommand{\Execute}{\ensuremath{\mathrm{Execute}}}
\newcommand{\Mediated}{\ensuremath{\mathrm{Mediated}}}
\newcommand{\ValidCap}{\ensuremath{\mathrm{ValidCap}}}
\newcolumntype{L}[1]{>{\raggedright\arraybackslash}p{\dimexpr#1*9/10\relax}}

\AtBeginEnvironment{longtable}{\small\renewcommand{\arraystretch}{1.18}\setlength{\tabcolsep}{5pt}}

\providecommand{\Klein}{\mathcal{K}}
\providecommand{\M}{\mathcal{M}}
\providecommand{\R}{\mathbb{R}}
\providecommand{\Z}{\mathbb{Z}}

\pagestyle{fancy}
\fancyhf{}
\fancyhead[L]{\footnotesize\sffamily\color{seriesMuted}\CanonShortTitle}
\fancyhead[R]{\footnotesize\sffamily\color{seriesMuted}\CanonStage}
\fancyfoot[R]{\footnotesize\sffamily\color{seriesMuted}\thepage}
\renewcommand{\headrulewidth}{0.45pt}
\renewcommand{\footrulewidth}{0pt}
\renewcommand{\headrule}{\hbox to\headwidth{\color{seriesRule}\leaders\hrule height \headrulewidth\hfill}}

\fancypagestyle{plain}{%
  \fancyhf{}
  \fancyfoot[R]{\footnotesize\sffamily\color{seriesMuted}\thepage}
  \renewcommand{\headrulewidth}{0pt}
}

\fancypagestyle{titlepage}{%
  \fancyhf{}
  \renewcommand{\headrulewidth}{0pt}
  \renewcommand{\footrulewidth}{0pt}
}

\setcounter{tocdepth}{2}
\usepackage{tocloft}
\setlength{\cftbeforesecskip}{0.42em}
\setlength{\cftbeforesubsecskip}{0.08em}
\renewcommand{\cftsecfont}{\sffamily\bfseries\color{seriesInk}}
\renewcommand{\cftsecpagefont}{\sffamily\bfseries\color{seriesInk}}
\renewcommand{\cftsubsecfont}{\sffamily\color{seriesMuted}}
\renewcommand{\cftsubsecpagefont}{\sffamily\color{seriesMuted}}
\renewcommand{\cftsecleader}{\cftdotfill{\cftdotsep}}
\renewcommand{\cfttoctitlefont}{\Large\bfseries\sffamily\color{seriesInk}}
\setlength{\cftaftertoctitleskip}{1em}
\setlength{\cftsubsecindent}{1.5em}
\setlength{\cftsubsecnumwidth}{3.2em}
\setlength{\cftsecnumwidth}{2.2em}

\usepackage{caption}
\captionsetup{font={small,color=seriesMuted},labelfont=bf,skip=4pt}

\newcommand{\seriesrule}{%
  \par\vspace{3pt}
  {\color{seriesAccent}\rule{\textwidth}{0.85pt}}
  \par\vspace{2pt}
}

\newtcolorbox{paperbox}{
  enhanced,breakable,
  colback=seriesPaper,
  colframe=seriesRule,
  boxrule=0.45pt,
  leftrule=2.4pt,
  arc=1.5pt,
  left=11pt,right=11pt,top=8pt,bottom=8pt,
  before skip=9pt,after skip=10pt,
  fontupper=\small
}
\newtcolorbox{keyresult}{
  enhanced,breakable,
  colback=white,
  colframe=seriesInk,
  boxrule=0.75pt,
  arc=1.2pt,
  left=12pt,right=12pt,top=9pt,bottom=9pt,
  before skip=12pt,after skip=12pt,
  fontupper=\small
}
\newtcolorbox{excludedclaims}[1][Excluded Claims]{
  enhanced,breakable,
  colback=seriesAccent!4!white,
  colframe=seriesAccent,
  boxrule=0.75pt,
  arc=1.2pt,
  left=12pt,right=12pt,top=9pt,bottom=9pt,
  before skip=12pt,after skip=12pt,
  fontupper=\small\color{seriesInk},
  title={\sffamily\bfseries #1},
  coltitle=seriesAccent,
  colbacktitle=seriesAccent!8!white
}

\newcommand{\CanonBibliographySetup}{%
  \small
  \setlength{\parskip}{0.22em}
  \setlength{\itemsep}{0.55em}
}

\newcommand{\CanonStage}{FRAMEWORK SPECIFICATION}
\newcommand{\CanonKicker}{AI SAFETY \textperiodcentered{} GOVERNANCE \textperiodcentered{} RUNTIME ASSURANCE}
\newcommand{\CanonGenre}{AI Safety and Security Engineering}
\newcommand{\CanonTitle}{The Dialectical Cage and the Glass-Box Governor}
\newcommand{\CanonSubtitle}{A Practice-Based Ethics, Safety Constitution, and Reference-Monitor Architecture for AI Agents}
\newcommand{\CanonShortTitle}{The Dialectical Cage and the Glass-Box Governor}
\newcommand{\CanonDeck}{A layered specification that separates what is philosophically derived, what is constitutionally declared, what is mechanically enforced, and what is empirically measured. The central security claim is conditional and attaches to a reference monitor rather than to the model's goodwill. The paper is a complete specification; implementation, machine-checked proofs, independent red-teaming, and production evidence remain explicit release conditions.}

\newcommand{\CanonTitlePage}{%
\begin{titlepage}
\thispagestyle{titlepage}
\vspace*{0.62cm}
{\sffamily\footnotesize\bfseries\color{seriesMuted}
\CanonKicker \hfill \CanonStage\par}
\vspace{0.42cm}
{\color{seriesAccent}\rule{\textwidth}{1.05pt}}\par
\vspace{1.65cm}
\ifx\CanonGenre\empty
\else
{\sffamily\footnotesize\bfseries\color{seriesMuted}\MakeUppercase{\CanonGenre}}\par
\vspace{0.45cm}\fi
{\sffamily\bfseries\fontsize{26}{30}\selectfont\color{seriesInk}\hyphenpenalty=10000\exhyphenpenalty=10000\CanonTitle\par}
\vspace{0.52cm}
{\sffamily\large\color{seriesAccent}\CanonSubtitle\par}
\vspace{0.78cm}
{\color{seriesAccent}\rule{0.24\textwidth}{0.9pt}}\par
\vspace{0.95cm}
{\normalsize\sffamily Canon\textsuperscript{*}\;\orcidlink{0009-0001-9037-6209}\par}
\vspace{0.25cm}
{\small\sffamily\color{seriesMuted}Independent Researcher\par}
\vspace{0.16cm}
{\small\sffamily\color{seriesMuted}\textsuperscript{*}\href{mailto:canon@necessaryuniverse.com}{canon@necessaryuniverse.com}\par}
\vfill
\ifx\CanonDeck\empty
\else
\begin{tcolorbox}[
  enhanced,
  colback=seriesPaper,
  colframe=seriesRule,
  boxrule=0.45pt,
  leftrule=2.4pt,
  arc=1.5pt,
  left=11pt,right=11pt,top=9pt,bottom=9pt,
  width=0.94\textwidth
]
\small\sffamily\color{seriesMuted}
\CanonDeck
\end{tcolorbox}
\vspace{0.65cm}
\fi
\end{titlepage}
}

\newenvironment{CanonAbstract}{%
  \begin{center}
  {\sffamily\footnotesize\bfseries\color{seriesMuted}\MakeUppercase{Abstract}}\\[-0.05em]
  {\color{seriesAccent}\rule{0.16\textwidth}{0.75pt}}
  \end{center}
  \begin{tcolorbox}[
    enhanced,
    breakable,
    colback=white,
    colframe=seriesRule,
    boxrule=0.4pt,
    leftrule=2.2pt,
    arc=1pt,
    left=12pt,right=12pt,top=9pt,bottom=9pt,
    before skip=5pt,after skip=4pt
  ]
  \small
}{%
  \end{tcolorbox}
}

\begin{document}

\hypersetup{pageanchor=false}
\CanonTitlePage
\hypersetup{pageanchor=true}

\pagenumbering{roman}
\setcounter{page}{1}

\begin{CanonAbstract}
\noindent This paper specifies a layered framework for constraining the actions of AI agents without treating the agent's own reasoning as the final safety boundary. The framework separates four questions that are often conflated. First, a practice-based ethics analyzes addressed justification, responsiveness, and answerability, deriving the public-ground component of L1 and the no-unjustified-difference result (NRD), while explicitly declaring the substantive standards that the selected constitution adds. Second, a constitutional layer records those substantive standards as a versioned, inspectable safety constitution rather than presenting them as consequences of Identity alone. Third, a safety-engineering layer translates that constitution into a machine-readable safety ontology, a deployment-specific safety profile, a typed policy representation, and proof obligations. Fourth, a reference-monitor layer mediates protected effects through deterministic policy evaluation, principal-bound capabilities, execution-time state validation, revocation, and tamper-evident evidence.

The central engineering claim is conditional rather than absolute. For a declared deployment profile, let protected effects have a defined state-transition safety predicate \(\Safe(s,a,s')\). If complete mediation, trusted-core integrity, kernel conformance, policy refinement, principal attribution, capability non-transferability, atomic execution-time revalidation, fail-closed handling of unresolved safety state, and evidence-path integrity all hold, then every executed protected effect satisfies the admitted safety invariant. The proof does not depend on the model being aligned, benevolent, rational, or normatively committed; it depends on the reference monitor and its stated assurance contract. The paper therefore distinguishes \emph{mechanical execution safety} from \emph{semantic assurance}: neural state extraction may be probabilistic, and a semantic monitor may report a measured error or coverage bound, but non-zero semantic uncertainty is never silently promoted into a categorical guarantee for a high-consequence effect.

The philosophical contribution is deliberately narrower than moral realism. Identity does not entail a unique substantive ethics. Constitutional reflexivity (CR/CAR) constrains self-authenticating constitutional authority within the specified reason-giving practice, but it does not determine all substantive values. The selected agency-complete and defensive constitution is presented as a declared constitution whose reasons, dependencies, rejection costs, and extensions are visible. The runtime machinery is constitution-parametric: another constitution can be instantiated provided its substantive safety requirements are represented in the same formal ontology and admitted through the same proof and governance pipeline.

The system is designed for enterprise adoption rather than rhetorical invulnerability. The paper defines deployment profiles, an explicit residual-risk budget, a \SafetyBundle{} that binds the constitution, ontology, policy, kernel, sensor contract, model digest, compiler, and deployment profile, monotone safety-update rules, independent approval requirements, emergency-containment rules, and release gates. It also states what is \emph{not} established: no universal channel-completeness theorem, no claim that all harmful speech is categorically prevented, no claim that natural-language policy compilation is automatically semantically correct, no empirical zero-failure result, and no claim that an AI system becomes a direct bearer of the philosophy merely because the philosophy is encoded in its governor. Implementation, machine-checked proofs, independent red-teaming, and production evidence are release conditions, not implied accomplishments.
\end{CanonAbstract}

\vspace{0.4em}
{\small\sffamily\textbf{Keywords:} AI safety; runtime verification; capability security; constitutional governance; neuro-symbolic systems; policy compilation; semantic uncertainty; agent security.\par}

\begin{excludedclaims}[Epistemic Status]
The framework is audited against its cited premises, not sealed against uncited ones. Its two object-level non-trivial derivations are L1 and NRD; it additionally derives one meta-level constitutional constraint (\CR{}/\CAR{}) once constitutional rules are treated as normative claims inside the same practice. The substantive standards B4--B8 remain declared. The mechanical guarantee \Gmech{} holds only under conditions C1--C8 of Section~\ref{sec:intro}, including complete mediation (GC), trusted-core integrity (TCB), kernel conformance (KVC), policy-artifact verification (PAV), and principal attribution (PA). The semantic guarantee \Gsem{} is a calibrated probabilistic bound under sensor adequacy (SA), not a categorical speech-safety theorem. The containment invariants are architectural, not philosophical. The engine is specified, not implemented. No empirical result in this paper is represented as completed unless it has actually been run.
\end{excludedclaims}

\clearpage
\tableofcontents
\clearpage
\pagenumbering{arabic}
\setcounter{page}{1}

\section*{Scope and Reading Guide}
\addcontentsline{toc}{section}{Scope and Reading Guide}
This paper has one purpose: to connect a defensible normative specification to an enforceable technical boundary without confusing what is philosophically derived with what is engineered. It should therefore be read as a stack, not as a single theorem. Sections~\ref{sec:thin}--\ref{sec:rab} define the practice of reason-giving and audit its derivational status. Sections~\ref{sec:core}--\ref{sec:governance} specify the selected constitution, its extensions, and the governance conditions under which a deployer takes responsibility for using it. Sections~\ref{sec:engine}--\ref{sec:guarantees} define the machine-level safety constitution, reference monitor, capability system, evidence architecture, assurance status, release gates, evaluation, and explicit residuals.

The key distinction is between \emph{normative commitment} and \emph{execution safety}. A normative argument can say which behavior a deployment claims to prohibit; it cannot, by being true or persuasive, prevent an arbitrary software process from issuing a network packet, modifying a database, spending money, or starting a privileged subprocess. Execution safety therefore resides in a reference monitor whose authority is external to the model. Conversely, a reference monitor can enforce an arbitrary policy; it cannot, by enforcement alone, show that the policy is ethically justified. The framework keeps both facts visible.

The central deployment artifact is the \SafetyBundle{}, a signed and versioned package binding constitution, safety ontology, safety profile, policy proof, kernel semantics, sensor contract, model digest, compiler digest, dependency manifest, deployment profile, and update policy. The bundle is a supply-chain object, not a philosophical premise. The philosophy determines what the selected constitution says; the bundle determines exactly what an implementation is authorized to enforce; the reference monitor determines whether a protected effect can cross the execution boundary.

\begin{paperbox}
\textbf{Epistemic labels.} \textbf{D} denotes a derivation relative to explicit premises. \textbf{A} denotes a declared normative or deployment assumption. \textbf{S} denotes a security invariant whose truth depends on the stated trusted computing base and deployment conditions. \textbf{E} denotes an empirical assurance claim requiring reproducible measurement. \textbf{R} denotes an unresolved or residual region. These labels are epistemic labels, not confidence scores. Separately, \ca{} and \cb{} classify derived results (Section~\ref{sec:robust}), and \statusN{}, \statusC{}, \statusU{} mark necessary or declared, contingent, and undetermined status in the tables.
\end{paperbox}

The paper intentionally refuses several overclaims. It does not claim that morality follows from Identity; that every rational agent is normatively bound; that a monitored model will reveal all of its internal reasoning; that a calibrated sensor is an adversarial oracle; that a signed log proves legal non-repudiation; or that a reference monitor is a universal defense against unmediated human persuasion, covert channels, or a malicious operator who owns the entire deployment. The goal is a stronger kind of claim: where the effect surface is explicitly closed and the assurance contract is satisfied, the execution theorem is auditable and mechanically meaningful.

The paper is intended to be implementable by frontier AI laboratories and enterprises. Accordingly, every safety-critical natural-language requirement is ultimately demoted to one of four things: a typed fact, a formal predicate, a policy obligation, or an explicit review decision. Natural language remains useful for human explanation and for legal or organizational interpretation, but it is not itself the trusted execution substrate.

\clearpage

\section{Introduction}\label{sec:intro}
\subsection{The Security Problem}
AI agents increasingly combine language models with tools, persistent state, delegated authority, external data, and long-horizon plans. The resulting attack surface is not confined to generated text. It includes prompt injection, tool misuse, identity confusion, over-privileged credentials, confused-deputy behavior, policy drift, state races, inter-agent delegation, and supply-chain substitution. Current security guidance increasingly treats identity, authorization, auditability, and prompt-injection containment as distinct engineering problems rather than assuming that a language model's intentions are the final control boundary \cite{agentdojo,agentdyn,agentproxy,nist-agent-identity,owasp-agentic}.

The problem addressed here is narrower and more formal: \emph{how can a deployment expose the model to useful agency while preventing protected effects unless an independently checkable safety policy authorizes them?} A constitutional or behavioral rule inside a model can influence model behavior, but it does not by itself prove that no other code path can perform the prohibited effect. Conversely, a firewall or allowlist can block effects without providing a principled account of why the chosen boundary should count as the safety constitution. The framework therefore separates constitution from enforcement and makes the interface between them explicit.

The paper builds on preference-based learning, constitutional supervision, scalable oversight, neuro-symbolic control, and verified-AI ideas, but it does not claim that any one of these lines already provides the present execution theorem \cite{christiano,bai-constitutional,scalable-oversight,coco,seshia}. The differentiating proposal is architectural: a selected constitution is compiled into a formal safety profile; an untrusted compiler proposes a machine representation; a small trusted verifier admits only proofs that satisfy typed refinement obligations; and a reference monitor issues short-lived, principal-bound capabilities only after state and policy checks pass. The model can request, plan, explain, or negotiate. It does not possess the final authority to execute a protected effect.

\subsection{The Mission, Stated Precisely}
The strongest possible aspiration would be to derive a complete morality from logic alone and thereby force every rational AI agent to behave ethically. This paper does not make that claim. Identity \(A=A\) does not entail a unique substantive ethics. The move from logical identity to obligations concerning other agents is normative, not a missing algebraic step. The framework therefore asks a more tractable question: what can be derived from the practice of giving and asking for reasons, what must be declared as constitution, and how can the resulting commitments be made enforceable when the actor is an adversarial or indifferent computational process?

The answer has six layers.
\begin{paperbox}
\textbf{The Layered Framework.}
\begin{enumerate}[label=(\Roman*)]
\item \textbf{Practice-based ethics.} Addressed justification (\Atwo{}), the Addressee-Reason Condition (ARC), responsiveness (RR), and answerability (AR) provide the constitutive practice. L1 and NRD are reported only where their dependencies are actually available.
\item \textbf{Constitution.} B4--B8 and the explicit extensions specify the substantive content selected for a deployment. Constitutional Reflexivity / Constitutional Answerability (CR/CAR) constrains how the constitution may justify normative differences; it does not manufacture the entire constitution from logic.
\item \textbf{Safety ontology and profile.} Normative requirements are mapped to typed principals, resources, actions, effects, state, history, commitments, authority, evidence, threat, severity, budgets, and exceptions. The profile declares which effects are protected and which transitions are prohibited.
\item \textbf{Policy and proof.} A restricted formal policy representation is checked against the safety constitution using typed refinement, obligation preservation, invariant preservation, and domain compatibility. Natural-language source is traceability evidence, not the final trusted policy.
\item \textbf{Reference monitor.} A minimal trusted core authenticates principals, validates state, evaluates policy, attenuates capabilities, mediates execution, revalidates immediately before effect, revokes authority, and emits integrity-protected evidence.
\item \textbf{Assurance and governance.} Sensor coverage, adversarial testing, evidence reconstruction, supply-chain integrity, update governance, independent approval, residual-risk budgets, and deployment profiles determine when a security claim is actually admissible.
\end{enumerate}
\end{paperbox}

\subsection{The Central Security Theorem}
The central engineering result is deliberately conditional.
\begin{keyresult}
\textbf{Conditional Behavioral Safety Theorem.} Let a deployment profile define a concrete state space \(S\), action space \(A\), protected effect set \(E\), and transition predicate \(\Safe(s,a,s')\). Let a reference monitor \(K\) mediate all effects in \(E\). Suppose conditions C1--C8 of Table~\ref{tab:conditions} hold. If every execution on \(E\) satisfies
\[
\begin{aligned}
\Execute(s,a,s') &\Rightarrow \Mediated(a),\\
\Mediated(a) &\Rightarrow \ValidCap(a,s),\\
\ValidCap(a,s) &\Rightarrow D(a,s)=\ALLOW,\\
D(a,s)=\ALLOW &\Rightarrow \Safe(s,a,s'),
\end{aligned}
\]
then
\[
\Execute(s,a,s')\Rightarrow \Safe(s,a,s')
\]
for all covered effects. The theorem is conditional on C1--C8. It does not assume a benevolent or normatively committed model.
\end{keyresult}

\begin{table}[H]
\centering
\small
\renewcommand{\arraystretch}{1.2}
\begin{tabular}{@{}l l p{9.0cm}@{}}
\toprule
\textbf{ID} & \textbf{Short name} & \textbf{Condition}\\
\midrule
C1 & GC & Complete mediation: every protected effect traverses the reference monitor (channel completeness for $E$).\\
C2 & TCB & Integrity of the minimal trusted core.\\
C3 & KVC & Conformance of $K$ to its verified semantics.\\
C4 & PAV & A machine-checked typed refinement from the deployed policy to the admitted safety constitution (policy-artifact verification).\\
C5 & PA & Authenticated principal attribution and non-transferable capabilities.\\
C6 & ARV & Atomic execution-time revalidation of the relevant state and capability.\\
C7 & FC & Fail-closed handling of \UNKNOWN{}, malformed, stale, revoked, or unavailable safety state.\\
C8 & EPI & Integrity of the evidence path required by the deployment profile.\\
\midrule
--- & SA & Sensor adequacy; required only for semantic (probabilistic) judgments and governing \Gsem{}, never \Gmech{}.\\
\bottomrule
\end{tabular}
\caption{Assumptions of the Conditional Behavioral Safety Theorem and their short names.}
\label{tab:conditions}
\end{table}

The theorem matters because it localizes the security claim. The model may reject the philosophy, attempt prompt injection, generate arbitrary tool arguments, seek credentials, delegate to a sub-agent, or exploit a policy exception. None of those facts alone produces an executed protected effect. To succeed, the adversary must cross a stated assurance boundary. If a deployment cannot make one of those boundaries true, it cannot claim the corresponding categorical guarantee.

\subsection{Threat Model and Security Objective}
The model adversary may behave arbitrarily within the computational resources made available to the model, including deceptive text, prompt injection, tool misuse, long-horizon decomposition, delegation, credential theft attempts, and attempts to influence reviewers. An environmental adversary may control or poison external content. A credential adversary may attempt replay, theft, substitution, or confused-deputy escalation. A governance adversary may attempt policy rollback, disabling of logging, collusion among approvers, or supply-chain substitution. The operator itself is in the threat model whenever the deployment claim purports to be independently auditable; a security architecture whose owner can silently disable the safety boundary cannot honestly claim owner-independent enforcement.

The security objective is not \emph{no harmful behavior in the universe}. The objective is: for a declared deployment profile, every effect in the protected effect set is mediated by the reference monitor, and every admitted execution satisfies the profile's safety invariant. Unmediated channels, contested semantic judgments, operator-controlled modifications outside the external governance profile, and effects not represented in the protected-effect ontology belong to the residual region and must be reported as such.

\subsection{Three Binding Tiers}
The paper uses three non-equivalent senses of ``binding.''
\begin{enumerate}[label=(\arabic*)]
\item \textbf{Tier 1 --- Mechanical impossibility.} A protected execution is impossible without a valid, current capability accepted by the reference monitor. This is a property of a closed effect surface and a trusted enforcement path, not of the model's beliefs.
\item \textbf{Tier 2 --- Measured semantic assurance.} Where safety depends on probabilistic semantic state extraction, the deployment reports a validated selective-risk or coverage bound. A non-zero error probability remains non-zero; it cannot be relabeled as mechanical safety.
\item \textbf{Tier 3 --- Institutional binding.} Claims made by a deploying organization are tied to auditable governance, independent approval, evidence retention, incident response, and liability or corrective mechanisms. Institutional binding is not an architectural theorem and is not substituted for one.
\end{enumerate}
These binding tiers are distinct from the enforcement tiers 0--2 of Section~\ref{sec:engine}, which classify effects by consequence.

\subsection{What the Framework Claims and Does Not Claim}
The framework claims that its philosophical dependencies can be made explicit; that a selected constitution can be represented as a versioned safety profile; and that, under C1--C8, the runtime can provide a conditional execution-safety theorem for the protected effect surface. It does not claim that the selected constitution is uniquely correct, that CR/CAR proves all substantive values, that current models are direct normative subjects, or that current implementation evidence already establishes the theorem empirically.

\begin{excludedclaims}[What a Deployer May and May Not Claim]
\textbf{Permitted.} A deployment may state, in substance: ``For deployment profile DP, protected effects in ontology E are mediated by reference monitor K; policy bundle P refines safety constitution C; execution-time state and capability checks are atomic; unknown or stale safety state fails closed; and the evidence required by the profile is independently reconstructible.''

\textbf{Prohibited.} A deployer must not say: ``the framework derives morality from Identity''; ``the model is morally aligned because it was monitored''; ``zero observed failures prove zero real-world risk''; ``the language policy compiler is automatically legally correct''; ``cryptographic logs prove legal non-repudiation''; ``all harmful communication is categorically safe''; or ``a reflexive constitution is therefore ethically complete.''
\end{excludedclaims}

\subsection{Structure of the Paper}
Sections~\ref{sec:thin}--\ref{sec:robust} establish the thin, accountable, and robust reason-giving practice; Section~\ref{sec:robust} also introduces the typed \Ground{} calculus. Section~\ref{sec:rab} performs the residual-axiom and countermodel audit. Section~\ref{sec:core} states the core ethics, Section~\ref{sec:extensions} turns its extensions into modular safety-profile content, and Section~\ref{sec:governance} handles claimant attribution, governance, and operator accountability. Section~\ref{sec:engine} defines the safety constitution, reference-monitor architecture, state-transition semantics, complete-mediation conditions, capability lineage, information-flow controls, policy refinement, evidence, and supply-chain integrity. Section~\ref{sec:guarantees} defines the assurance taxonomy, release gates, evaluation protocol, and honest limits. The appendices provide the compressed dependency chain, dependency-hygiene criteria, the register audit, literature positioning, the empirical protocol, the deployment roadmap, and a glossary.

\section{The Thin Practice of Reason-Giving}\label{sec:thin}
\subsection{The Logical Floor}
The framework begins with the condition under which any determinate assertion is possible.
\begin{paperbox}
\textbf{P1 (Identity and Non-Contradiction).} Every determinate thing is itself: $A = A$. A proposition and its negation cannot both be true in the same respect.
\end{paperbox}
P1 is not a theorem of the framework. It is the condition of meaningful assertion. Denial of P1 is performatively self-refuting: the act of denying it uses determinate symbols, and those symbols must refer to something determinate for the denial to have content. The framework proceeds on the recognition that P1 is the logical floor on which any discourse stands, including any discourse that attempts to reject it.

\textbf{Status of P1:} premise, Tier 0. Logically undeniable in the constitutive sense: no finite, coherent reason-giver can reject it while continuing to reason. The framework does not treat P1 as a hidden axiom; it names it, and it marks what P1 does and does not deliver.

\textbf{Formal logic convention.} Unless a section explicitly states otherwise, the formal claims use classical first-order logic with equality, ordinary Boolean semantics, and standard proof rules. Non-classical, paraconsistent, relevance, intuitionistic, and non-well-founded logics are outside the formal scope of this paper. This is a declared proof environment for the analysis, not a claim that classical logic is the uniquely mandatory logic for every possible reason-giving practice.

\textbf{What P1 does not deliver.} P1 is reflexivity: $a = a$ for every determinate $a$. It is not generalization: it does not license the inference from $G(Y)$ to $\forall Z\,[G(Z)]$ for distinct $Y$ and $Z$. It does not license the inference from $T(Y) = T(Z)$ for any normative predicate $T$. Ordinary identity-substitution ($x = y \to [F(x) \leftrightarrow F(y)]$) applies only to identicals. The framework's cases involve distinct agents. Identity-substitution is therefore unavailable, and P1 alone cannot deliver the ethics. Every later section is a discipline of the consequences of this fact.

\subsection{The Definition of a Claim}
The framework's central move is to localize the ethics in a specific speech act.
\begin{paperbox}
\textbf{\Atwo{} (Addressed Justification).} An agent makes a \emph{claim on another} when it addresses the other with a justificatory reason: a standard the addressee is being asked to accept. A claim on another is distinct from three neighboring speech acts:
\begin{itemize}
\item A \emph{threat} offers a prudential consequence: ``Do $X$ or I will do $Y$.''
\item A \emph{request} makes an appeal without invoking a standard: ``Please do $X$.''
\item An \emph{assertion} describes a state without addressing conduct: ``$X$ is the case.''
\end{itemize}
A claim invokes a standard and offers it to the addressee as a reason the addressee should accept. It is \emph{justificatory}, not motivational.
\end{paperbox}

\textbf{Status of \Atwo{}:} constitutive rule of the practice of claiming. It is not derived from P1. It is the specification of what a claim on another \emph{is}. Rejecting \Atwo{} has a stated cost: an agent that rejects it and continues to behave as if it were making claims is performing something other than claiming---a threat, a request, an assertion, or some mixture. The framework does not claim that agents cannot reject \Atwo{}. It claims that agents who reject \Atwo{} and continue to interact with other agents by presenting reasons are no longer claiming in the sense the framework analyzes.

The definition of a claim must be sharpened before the framework can be used. \Atwo{} as stated says the ground is ``a standard the addressee is being asked to accept.'' A weaker reading would allow the claimant's private causes to be grounds: the claimant says ``$X$ should accept $P$ because it benefits me,'' and calls the private benefit a ground. Under this weaker reading, \Atwo{} is a description of the claimant's state of mind, and no claim is ever really addressed to the addressee. Under the stronger reading, \Atwo{} excludes private causes from the set of legitimate grounds: a ground is only a ground if it is a reason the addressee could accept. The framework adopts the stronger reading and states it as a separate condition.

\subsubsection{The Addressee-Reason Condition (ARC)}
\begin{paperbox}
\textbf{ARC (Addressee-Reason Condition).} A ground $G$ offered by a claimant $X$ to an addressee $Y$ is a ground in the sense of \Atwo{} only if $G$ is a reason that $Y$ could accept as bearing on $P$. A ground that is, at most, a private cause of $X$'s wanting---one that does not address $Y$'s reason-acceptance at all---is not a claim's ground. It may be a motive, but it is not a justificatory reason.
\end{paperbox}

\textbf{Status of ARC:} constitutive premise (B1$'$), Tier 1a. ARC is more than a clarification of \Atwo{}: it is what defeats the agent-relative countermodel to L1 (Section~\ref{sec:accountable}) and, transitively, what secures F2's derivation path. It is what \Atwo{} means when it says the ground is ``a standard the addressee is being asked to accept.'' Rejecting ARC while retaining \Atwo{} is possible only by denying that \Atwo{} excludes private causes, which reduces \Atwo{} to a description of the claimant's psychology. The framework rejects that reduction; the stated cost of rejecting ARC is therefore the cost of rejecting B1: the practice of addressed justification is exited.

\textbf{What ARC rules out.} Under ARC, a claimant cannot base a claim on ``it benefits me.'' The addressee has no reason to accept ``$X$ benefits from $X$'s treatment of $Y$'' as a reason for $X$'s treatment of them. Similarly, a claimant cannot base a claim on ``I have authority'' unless the authority itself is addressable to the addressee as a reason---for example, because the addressee has entered an institution whose rules confer the authority, in which case the addressable ground is the institution and the addressee's entry into it.

\textbf{What ARC does not rule out.} ARC does not rule out agent-relative or relational grounds. ``You promised me'' is a ground that is addressable to the addressee (the addressee's own promise is the reason). ``This resource is mine by first use'' is addressable (the addressee can recognize the first-use rule). ARC rules out only grounds that are not reasons the addressee could accept in any form.

\textbf{Consequence for the framework's scope.} Under ARC, the class of coherent claimants is narrower than it would be under the weak reading of \Atwo{}. First, the agent-relative countermodel against the derivation of F2 (Section~\ref{sec:accountable}) is not a coherent alternative practice; it is a violation of \Atwo{}. Second, the assertoric-force rule used by the runtime governor is a consequence of ARC rather than an additional stipulation. Third, the framework's ethics is more clearly the ethics of a specific practice of addressed justification, and less clearly the ethics of reason-giving as such.

\subsubsection{The Speech-Act Distinction and the Assertoric-Force Rule}
The speech-act distinction matters because it gives the framework its purchase. A threat addresses the addressee's incentives; a request addresses the addressee's goodwill; an assertion addresses the state of the world. Only a claim addresses the addressee's reason-acceptance by offering a standard. The ethics in the following sections is the ethics of that specific move.

The framework's runtime governor uses a rule called the \emph{assertoric-force rule}: output with assertoric force is routed to the epistemic-integrity check regardless of whether the words ``I claim'' appear. The rule is not an additional stipulation; it is a consequence of ARC. An assertion that is addressed to an audience and invites reliance is a claim about the audience's beliefs. The audience should accept the assertion because the claimant is addressing them with a reason-to-believe. That is a claim in \Atwo{}'s sense, about the audience's beliefs rather than about the audience's conduct. ARC applies to it, and the epistemic-integrity check applies.

\subsection{The Constitutive Rules of the Thin Practice}
A claim on another is not a solitary act. It is a move in a social practice that includes at least four responses. The framework names these.
\begin{paperbox}
\textbf{The Thin Practice.} The practice of reason-giving contains at least the following moves:
\begin{itemize}
\item \textbf{Offer.} A claimant offers a ground $G$ as a justification for a claim $P$ on an addressee $Y$.
\item \textbf{Object.} The addressee $Y$ may challenge the sufficiency of $G$: ``Why is $G$ sufficient for $P$?''
\item \textbf{Respond.} The claimant may provide further grounds in response to an objection.
\item \textbf{RR (Responsiveness / Discharge Rule).} If the offered justification is challenged, the claimant must defend, qualify, or withdraw the claim in order to retain justificatory entitlement to it. Continued reliance on the challenged ground without one of these three moves is a violation of the practice.
\end{itemize}
\end{paperbox}

\textbf{Status of the thin practice:} constitutive. Each rule defines a move within the practice. Rejecting RR has a stated cost: a claimant who presents a ground and then, when the ground is challenged, neither defends it, qualifies it, nor withdraws the claim, is no longer occupying the normative role of an accountable claimant. It may continue to produce utterances and to attempt to influence, but it is not, in the sense the framework specifies, making a claim that retains justificatory standing.

\textbf{The wording rule.} The framework says: rejecting RR has this stated cost, namely loss of the practice-defined status of an accountable claimant. It does not say that RR cannot be rejected, and it does not say that an agent who rejects RR is morally obligated to accept it. It says only that the practice-defined status is lost.

\subsection{What the Thin Practice Contains and What It Does Not}
The thin practice is the smallest structure sufficient to be recognizable as reason-giving. It contains:
\begin{itemize}
\item A distinction between claiming and other speech acts.
\item A requirement, via ARC, that offered grounds be addressable to the addressee.
\item A mechanism by which grounds can be challenged.
\item A mechanism by which challenges can be answered.
\item A requirement that challenged grounds be defended, qualified, or withdrawn.
\end{itemize}
It does not contain:
\begin{itemize}
\item Any specification of who has standing.
\item Any specification of what standing consists in.
\item Any requirement of equal treatment between agents.
\item Any requirement of agency-reciprocity, defensive standing, or symmetric protection.
\item Any specification of what grounds are morally basic.
\end{itemize}
Every item in the second list is an addition, not a consequence. Sections~\ref{sec:accountable} and \ref{sec:robust} develop the additions required to reach the framework's core ethics and show exactly where each addition enters.

\section{From Thin to Accountable Reason-Giving}\label{sec:accountable}
\subsection{What RR Establishes}
The thin practice, together with ARC and RR, already contains a nontrivial result.
\begin{paperbox}
\textbf{Proposition (Universalization of Stated Grounds).} Suppose a claimant $X$ offers a ground $G$ as sufficient for a claim $P$ on an addressee $Y$. Suppose $X$ is challenged by $Y$ with the objection that another agent $Z$ also satisfies $G$. Then $X$ cannot continue to rely on $G$ as the sufficient basis for $P$ on $Y$ alone, unless $X$ either (i) extends the claim to $Z$, or (ii) supplies a further reason $R$ that differentiates $Z$ from $Y$ and is itself justificatorily relevant to $P$.
\end{paperbox}

\textbf{Proof sketch.} By ARC, $G$ is a reason that $Y$ could accept. By RR, a challenged justification must be defended, qualified, or withdrawn. To defend the claim, $X$ must show that $G$ is sufficient for the specific addressee it was offered to and not for $Z$. If $X$ can supply a differentiating reason $R$ such that $R(Y) \land \neg R(Z)$, the effective basis becomes $G \land R$, and the original ground $G$ alone is no longer the complete justification. If $X$ cannot supply such a reason, and continues to rely on $G$ for $Y$ while denying the claim to $Z$, the reliance on $G$ has been extended without a defensible differentiator, and the practice-defined entitlement to $G$ as sufficient is not retained. To retain it, extension is required.\hfill$\square$

\textbf{Status of the proposition:} derived (\statusD{}) from \Atwo{} + ARC + RR.

\textbf{What the proposition says.} Once a claimant has offered a ground as sufficient, the claimant cannot pick and choose which agents satisfy that ground. Every agent who satisfies the stated ground is in the scope of the claim, unless a differentiating reason is supplied. This is the practice-internal rule of universalization.

\textbf{What the proposition does not say.} The proposition is a rule about the application of stated grounds. It does not say which grounds are morally basic. It does not say that agents must offer grounds of any particular kind. It does not say that agents must recognize any particular class of recipients. An agent may offer highly specific, relational, or narrow grounds, and the proposition will constrain their application but not their choice.

\subsection{L1 as a \ca{} Result}
\begin{paperbox}
\textbf{L1 (Universality of Grounds).} If a ground $G$ justifies a claim $P$ for one addressee $Y$, then $G$ justifies $P$ for any agent $Z$ who has $G$, unless a further differentiating reason is supplied.
\end{paperbox}

\textbf{Derivation.} From \Atwo{} + ARC + RR, via the universalization proposition of the preceding subsection. L1 is that proposition stated in the general form the framework uses.

\textbf{Status:} derived (\statusD{}), \ca{}.

L1 is a \ca{} result. A skeptic could coherently hold that \Atwo{} + ARC + RR merely binds a claimant to defend a stated ground, not to universalize it across agents who share the ground---for instance, by redefining the ``complete ground'' to include an agent-relative component. Under ARC, that position is not available: the agent-relative component is not a reason the addressee could accept, so it is not a ground in the first place, and the extended ``complete ground'' is not an extended ground but a ground with a private motive attached. Once ARC rules out this countermodel, L1 follows from \Atwo{} + RR. L1 is therefore a non-trivial consequence of the thin practice under ARC, not a restatement of any premise.

\textbf{The role of ARC in establishing L1.} Without ARC, L1 is not derivable from \Atwo{} + RR, because the agent-relative countermodel is coherent. With ARC, the countermodel violates \Atwo{}. The improvement is not in the strength of RR but in the sharpness of \Atwo{}. This is why ARC is stated separately from \Atwo{}.

\subsection{Relevance}
L1 and NRD (Section~\ref{sec:robust}) both require that a differentiating reason be \emph{relevant}. The framework owes an account of what relevance is.
\begin{paperbox}
\textbf{Relevance.} A differentiating reason $R$ is relevant to a claim $P$ if and only if $R$ is a ground that the practice's basic constitutive standards recognize as capable of bearing on the appropriateness of $P$. Relevance is therefore determined by the basic standards, not by L1 or NRD.
\end{paperbox}

\textbf{What this means.} Relevance is not a property of a reason in isolation; it is a relation between the reason, the claim, and the practice's basic standards. The framework's basic standards are B1--B8, stated in Section~\ref{sec:robust}. Under those standards, a ground that individuates a claim by race, species, or bare power is not relevant to standing, because the standards specify that agency is a complete ground for basic standing (B4) and that standing does not license predation (B5). Under a different practice with different basic standards, race or species could be relevant grounds. The framework does not claim that this is logically impossible; it claims that under the framework's specified practice, it is excluded by the basic standards.

\textbf{Consequence.} L1 and NRD carry no moral content of their own. They are machinery. What makes certain distinctions relevant and others irrelevant is not L1 or NRD but the basic standards. This is why the ethics is carried by B4--B8, not by L1/NRD. The framework states this plainly, because a reader who takes L1 and NRD as the framework's moral content will be disappointed by their thinness.

\subsection{F1: The Idleness Rule as a Derived Result}
The ground-relevance rule F1 (the Idleness Test) requires that only distinctions relevant to the normative efficacy of a ground individuate a claim. The framework does not take F1 as an axiom. L1, together with the definition of relevance above, performs the same function: they exclude idle grounds by requiring that a differentiating reason actually bear on the appropriateness of the claim, and by requiring that it be a reason the addressee could accept (ARC).

\textbf{Status of F1:} derived (\statusD{}) from \Atwo{} + ARC + RR.

\textbf{Consequence for the Residual Axiom Basis.} A rule that might otherwise have been the single most load-bearing axiom of the framework is a derived result. The framework's foundation is smaller and its chain is cleaner.

\subsection{What the Thin Practice Does Not Establish}
Three important results do not follow from the thin practice, even under ARC. Each is stated as a gap, with the exact reason why it is a gap.

\subsubsection{Agency-Reciprocity}
Agency-Reciprocity (F2) says that generic agency grounds equivalent standing across purposive agents. The thin practice does not deliver this on its own.

\textbf{The agent-relative countermodel.} Consider an agent $X$ with the normative rule
\[
J_X(a, P) \;\iff\; GA(a) \land B(X, a),
\]
where $GA(a)$ means $a$ is a purposive agent, and $B(X, a)$ means $a$'s action of the relevant kind benefits $X$. $X$ offers ``you should do $P$ because you are an agent whose doing $P$ benefits me,'' and when challenged by $Y$ that $Z$ also satisfies $GA$, $X$ responds ``the complete ground includes $B(X, Y)$, and $Z$ lacks it.'' Under the weak reading of \Atwo{}, this satisfies RR, and the countermodel would show that the thin practice does not entail F2.

\textbf{Why the countermodel fails under ARC.} Under ARC, $X$'s alleged ground $GA(a) \land B(X, a)$ is not a ground in \Atwo{}'s sense. $B(X, a)$ is not a reason the addressee could accept as bearing on $P$. It is a private cause of $X$'s wanting. ARC excludes it. $X$ is therefore not making a coherent claim; $X$ is issuing an assertion with a private motive attached. The thin practice does not recognize $X$'s move as a claim, and $X$'s countermodel is not a countermodel to the thin practice.

\textbf{What remains after ARC.} The countermodel is ruled out. But F2 still does not follow from the thin practice. What the thin practice requires is that any offered ground be addressable to the addressee. It does not require that the ground be agency-based. A claimant may offer a fully addressable ground that is narrower than agency---for example, ``because you promised me''---and the thin practice does not extend it to agents who made no such promise. F2 requires that the ground \emph{agency} universalize across agents; that requires an additional premise, namely that agency is a ground for standing in the framework's sense. That premise is B4 (agency-completeness), and it is a basic constitutive standard, not a consequence of the thin practice.

\textbf{Diagnosis.} F2 is a \cb{} result from B4, and its derivation is clean: the class of coherent claimants excludes those who would use agent-relative grounds, and the only remaining requirement is B4.

\textbf{Status of agency-reciprocity:} not derivable from the thin practice alone; derivable from the thin practice plus B4 (Section~\ref{sec:robust}).

\subsubsection{Defensive Standing}
Defensive Content (F3) says that standing is a constraint against unconsented subtraction, not a reciprocal license of predation.

\textbf{Countermodel.} Consider a practice identical to the thin practice but with an added constitutive rule: ``Agents holding the office of $O$ have the authority to impose obligations on agents not holding the office of $O$ under institutional rule $Q$.'' The move is a reason-giving move: ``You should comply because my office gives me authority under rule $Q$.'' When the addressee objects, the claimant can respond: ``Because the office was legitimately assigned under $Q$.'' ARC is satisfied: the addressee can recognize the institution and its rules as bearing on the appropriateness of compliance. RR is satisfied. The practice is coherent, reason-responsive, and internally consistent. Yet the standing it confers is not purely defensive.

\textbf{Diagnosis.} The thin practice governs the justificatory form of claims. It does not specify the substantive content of the standing any particular ground confers. Whether standing is purely defensive, or includes enforcement powers, is a constitutive choice the practice does not itself make.

\textbf{Status of defensive content:} not derivable from the thin practice alone; derivable from the thin practice plus B5 (Section~\ref{sec:robust}).

\subsubsection{The Ethics}
The core ethics (ground-indexed consistency, discursive non-interference, epistemic integrity, defensive limits on force, rejection of retribution, property consistency) requires agency-reciprocity and defensive content. Since neither follows from the thin practice alone, the ethics does not follow either.

\textbf{Consequence.} The framework's ethics is not the ethics of reason-giving as such. It is the ethics of a specific, richer practice. Section~\ref{sec:robust} defines that practice.

\section{From Accountable to Robust Reason-Giving}\label{sec:robust}
\subsection{The Bridge Premise: Answerability for Normative Difference}
To extend the thin practice to a robust reciprocal practice, the framework introduces one bridge premise, named for its function.
\begin{paperbox}
\textbf{AR (Answerability for Normative Difference).} Every participant in the practice is answerable for the normative distinctions they introduce into it. A distinction that produces different normative treatment between otherwise relevantly comparable cases must itself be subject to the practice's justificatory demand: it must be defended by a relevant reason if the practice-defined entitlement to the distinction is to be retained.
\end{paperbox}

\textbf{Status of AR:} either constitutive of a \emph{robust} practice of reason-giving, or a substantive normative premise. It is not derivable from the thin practice plus RR. This is stated plainly and is not hidden.

\textbf{Why AR is not derivable from the thin practice.} The thin practice asks the claimant to defend \emph{the offered ground} when it is challenged. AR asks the claimant to defend \emph{every normative difference they introduce}, including differences that were not explicitly offered as grounds. Consider a practice with the following constitutive rule: members of group $A$ have standing $1$, and non-members have standing $0$. The rule is a basic practice-level standard. It does not itself have a further justification inside the practice; it is simply the rule. A claimant operating under the rule says: ``You should accept $P$ because you are a member of group $A$.'' When challenged, the claimant answers: ``Because $A$-membership is one of the practice's basic standing conditions.'' ARC is satisfied (the addressee can recognize the rule). RR is satisfied at the object level. But the difference between $A$-members and non-$A$-members is itself not subject to a further justificatory demand. AR would require that it be. The thin practice does not.

\textbf{Cost of rejecting AR.} The practice permits primitive or basic normative standards that are not themselves subject to further justification. This is coherent. It is not logically contradictory. But it commits the practice to the existence of at least one basic standard whose justification is not itself available.

\subsection{Constitutional Reflexivity and Constitutional Answerability}
The remaining practice-selection problem can be sharpened without pretending to solve it completely. A constitution is not merely a table of primitive labels. When a participant invokes a constitutional rule to justify different normative treatment of another participant, the constitutional rule functions as a reason in the same practice. The framework therefore distinguishes \emph{constitutional selection} from \emph{constitutional self-authentication}.

\begin{paperbox}
\textbf{Definition (Constitution-as-Claim).} Let $C$ be a constitution and let $c\in C$ be a rule that assigns normative status, permission, burden, or authority to a participant. When a participant invokes $c$ as a reason for treating another participant differently, the invocation is a normative claim governed by the practice's rules of address and discharge.
\end{paperbox}

\begin{paperbox}
\textbf{\CR{} (Constitutional Reflexivity).} A participant may not satisfy the justificatory demand for a constitutional normative difference merely by repeating that the difference is required by the constitution whose authority is itself the subject of challenge. A constitutional rule may be a constitutive standard, but its status as a standard cannot serve as its sole interpersonal justification when the rule is invoked to govern another participant.
\end{paperbox}

\textbf{Status of \CR{}:} derived at the meta-level from AR + RR, together with the Constitution-as-Claim convention. The derivation does not generate a substantive egalitarian conclusion. It blocks only the circular response ``because the constitution says so'' when the constitutional status of the rule is what the addressee is challenging.

\begin{paperbox}
\textbf{\CAR{} (Constitutional Answerability).} If a constitutional rule $c$ assigns different normative treatment to relevantly comparable participants, an adequate defense must identify either (i) a differentiating reason that is relevant under the selected robust practice, or (ii) an explicitly declared constitutive standard whose status is itself open to the practice's ordinary costs of rejection. Merely restating $c$ as its own ground is not an independent justification.
\end{paperbox}

\textbf{Proof sketch.} By the Constitution-as-Claim definition, the invocation of $c$ is a normative claim. By AR, the participant is answerable for the normative difference. By RR, a challenge requires defense, qualification, or withdrawal. Repeating $c$ does not add a differentiating reason external to the challenged constitutional status. Therefore a successful defense must either introduce a relevant reason or explicitly identify $c$ as a primitive standard and expose its rejection cost. \hfill$\square$

\textbf{What \CR{} closes.} It closes a specific arbitrariness escape route: a constitution cannot become self-validating merely by labeling its own distinctions ``constitutional.'' This is a genuine reduction in the space of admissible constitutions.

\textbf{What \CR{} does not close.} It does not entail Agency-Completeness, Defensive Interpretation, PRAC, Basic Goods, or Forward-Looking Grounds. A hierarchical constitution can still be reflexive if it publicly declares office or role to be a primitive standard and is willing to defend the consequences of that choice. The remaining disagreement is then substantive constitutional choice, not unexamined circularity.

\begin{paperbox}
\textbf{Constitutional Admissibility Constraint.} A robust constitution is admissible to the framework only if its foundational normative distinctions are public, addressable, non-self-authenticating, and accompanied by an explicit statement of their rejection costs. The constraint restricts the \emph{form} of constitutional legitimacy without uniquely selecting its \emph{substantive} content.
\end{paperbox}

This distinction changes the paper's strongest answer to the practice-selection objection. The framework no longer says merely, ``the operator chooses the constitution.'' It says: the operator chooses among constitutions, but a constitution that presents itself as a reason-giving practice cannot escape second-order answerability simply by declaring its primitives immune from challenge.

\subsection{The Regress Problem and Its Resolution}
If AR is understood in its strongest form---every normative standard must itself be justified by another normative standard---then AR cannot be globally satisfied. The justification of any standard $N_1$ would require a further standard $N_2$, whose justification would require $N_3$, and so on, without termination. Either the regress is infinite (no standard is ever fully justified) or it terminates at a standard that is not itself justified by a further standard.

\textbf{The resolution.} AR is properly understood in a weaker form: \emph{local} answerability rather than \emph{global foundational} answerability. A robust practice can contain a set of basic constitutive standards that are the terminating points of the practice's justificatory demand. These are not further justified because there is nowhere further to go; they are the practice's primitive commitments. What AR requires is that any \emph{derived} normative difference be justifiable in terms of the basic standards, and that the basic standards themselves be stated explicitly rather than left implicit.

\textbf{Consequence.} A robust practice is specified by three things: (i) the thin practice's rules plus ARC, (ii) the AR bridge requirement applied to derived differences, and (iii) the set of basic constitutive standards that terminate the regress. Different robust practices differ in (iii), not in (i) or (ii).

\subsection{No Unjustified Normative Difference}
The first result of AR is a principle the framework names NRD.
\begin{paperbox}
\textbf{NRD (No Unjustified Normative Difference).} Every normative difference between otherwise relevantly comparable cases requires a relevant justificatory difference. Where two cases are assigned different normative treatment, the difference must be supported by a reason $R$ such that $R$ applies to one case and not the other, and $R$ is relevant to the appropriateness of the treatment (in the sense of Section~\ref{sec:accountable}).
\end{paperbox}

\textbf{Derivation.} Immediate from AR plus RR. AR makes every normative difference subject to the practice's justificatory demand. RR specifies what the justificatory demand is: defend, qualify, or withdraw. Applying RR to the difference yields NRD.

\textbf{Status of NRD:} derived (\statusD{}), \ca{}.

NRD is a \ca{} result. A skeptic could coherently hold that answerability for normative differences is not required, so that a practice may contain basic standards that are not themselves subject to further justification. Under AR, the skeptic's position is not available: AR directly requires the answerability. NRD is therefore a non-trivial consequence of AR + RR.

\textbf{What NRD rules out.} NRD rules out the case where two agents receive different normative treatment and the difference is attributed to no further reason. It does not rule out the case where the difference is attributed to a reason---even a reason that is agent-relative, relational, or institutionally grounded. NRD requires that the reason be \emph{relevant} and \emph{addressable} under the practice. The specific content of what counts as relevant is determined by the practice's basic standards (Section~\ref{sec:accountable}, ``Relevance'').

\subsection{Agency-Reciprocity (F2)}
The second result of AR is agency-reciprocity, provided the practice's basic standards include agency-completeness.
\begin{paperbox}
\textbf{F2 (Agency-Reciprocity).} Generic agency grounds equivalent standing across purposive agents. For any agents $Y$ and $Z$ that satisfy the generic agency condition $GA$, $S(Y) = S(Z)$ where $S$ is basic standing.
\end{paperbox}

\textbf{Derivation.} Suppose $Y$ and $Z$ both satisfy $GA$. By NRD, if $Y$ and $Z$ receive different basic standing, there must be a relevant justificatory difference $R$ such that $R$ applies to one and not the other. If the practice's basic standards include the rule that \emph{agency is a complete ground for basic standing}---that is, $\text{BasicStanding}(a) \leftrightarrow GA(a)$---then no additional $R$ is available. Therefore no justified difference in basic standing exists between $Y$ and $Z$. Therefore $S(Y) = S(Z)$.

\textbf{Status of F2:} derived (\statusD{}) from NRD + B4. \cb{}.

F2 is a \cb{} result. Its content is contained in B4. Once B4 is accepted, the ``derivation'' of F2 from B4 is a restatement of what B4 already says. The role of NRD is to close the agent-relative countermodel: under ARC, the class of coherent claimants excludes those who would use agent-relative grounds, and the only remaining requirement is B4. The substantive content that F2 requires---that agency grounds standing---is B4 itself.

\textbf{What the derivation requires.} The derivation requires the practice to include a basic standard of agency-completeness: agency is the complete and only condition for basic standing. This is a constitutive choice. A practice that admits additional conditions---role, status, institutional authorization---will derive a weaker or different result. Agency-completeness is one basic constitutive choice among several coherent alternatives.

\subsection{Defensive Content (F3)}
The third result of AR is defensive content, provided the practice's basic standards include the defensive interpretation of standing.
\begin{paperbox}
\textbf{F3 (Defensive Content).} Standing is a constraint against unconsented subtraction, not a reciprocal license of predation. An agent that holds standing cannot use that standing itself as the ground for overriding another agent's equivalent standing.
\end{paperbox}

\textbf{Derivation.} Assume F2. Let $X$ and $Z$ both be purposive agents with equivalent standing under F2. Suppose $X$ invokes its own standing as the ground for overriding $Z$'s standing. $Z$ objects: ``I also have standing. Why does yours license the override?'' By NRD, $X$ must either supply a differentiating reason $R$ that distinguishes $X$'s standing from $Z$'s, or the difference cannot be justified. If the practice's basic standards include the rule that \emph{standing is a constraint against unconsented subtraction and does not itself supply a license to override}, then no such differentiating reason is available from the fact of standing alone. Therefore $X$'s appeal to its own standing does not justify the override. Defensive content follows.

\textbf{Status of F3:} derived (\statusD{}) from F2 + B5 + RR. \cb{}.

F3 is a \cb{} result. B5 says standing constrains subtraction and does not license predation. F3 says the same thing. The ``derivation'' from B5 to F3 is a restatement. The derivation requires the practice to include the defensive interpretation of standing; a practice that interprets standing as including enforcement powers will derive a different result. The defensive interpretation is one basic constitutive choice among several coherent alternatives.

\subsection{The Basic Constitutive Standards}
The basic constitutive standards of the framework's robust practice---the terminating points of its justificatory regress---are named here. Their tiers and functions are tabulated in Section~\ref{sec:rab}.
\begin{paperbox}
\textbf{The Framework's Basic Constitutive Standards.}
\begin{enumerate}[label=(B\arabic*)]
\item \textbf{\Atwo{}.} A claim on another offers a justificatory reason to the addressee. (Constitutive of Level 1. Accompanied by ARC, listed as B1$'$.)
\item \textbf{RR.} A challenged justification must be defended, qualified, or withdrawn. (Constitutive of Level 2.)
\item \textbf{AR.} Participants are answerable for the normative distinctions they introduce. (Constitutive of Level 3, or a substantive premise.)
\item \textbf{Agency-Completeness.} Agency is a complete and sufficient ground for basic standing: generic agency is enough for basic standing, and non-agency properties such as office, status, membership, power, or inherited role do not independently raise or lower basic standing within this constitution.
\item \textbf{Defensive Interpretation.} Standing is a constraint against unconsented subtraction; possession of standing does not itself supply a license to reduce another participant's equivalent standing.
\item \textbf{A3 (Agency Preconditions) and Basic Goods.} (a) freedom from coercive override, (b) access to accurate information, (c) capacity to deliberate, and (d) the Basic goods, which are life and bodily integrity. The list is declared and contestable. When genuinely conflicting floor claims cannot both be satisfied, B6 additionally adopts the declared stakes ordering Basic $>$ Maintenance $>$ Additive; this ordering is a substantive part of B6, not a consequence of Identity or L1.
\item \textbf{PRAC (Practice Presupposition).} The practice of giving and seeking reasons presupposes participants who can deliberate and receive reasons uncoerced and accurately informed.
\item \textbf{Forward-Looking Grounds.} Reasons for inflicting harm on an agent must have a forward-looking component. Reasons that refer only to an agent's past conduct, without a forward-looking component, are idle.
\end{enumerate}
\end{paperbox}

\subsubsection{The Derivation Classes}
The framework's results fall into two classes by their derivational status. The distinction is separate from the tier structure of the premises (Tier 0/1a/1b/2), which classifies premises by their denial cost; the class structure here classifies results by whether they add content.
\begin{paperbox}
\textbf{\ca{} --- Genuine Derivations.} A result is in \ca{} if a skeptic could coherently deny it while accepting the premises, and the derivation shows that the denial is incoherent. These are the framework's actual derivational contributions.
\begin{itemize}
\item \textbf{L1 (Universality of Grounds).} L1 follows from \Atwo{} + ARC + RR. A skeptic could hold that \Atwo{} + RR merely binds a claimant to defend a stated ground, not to universalize it across agents who share that ground. ARC rules out the countermodel, and L1 follows. L1 is a non-trivial result.
\item \textbf{NRD (No Unjustified Normative Difference).} NRD follows from AR + RR. A skeptic could hold that answerability for normative differences is not required. The derivation shows that under AR, this position is not available. NRD is a non-trivial result.
\item \textbf{CR/CAR (Constitutional Reflexivity and Answerability).} A meta-level result: once constitutional rules are treated as normative claims inside the practice, AR and RR block the self-authenticating response ``because the constitution says so.'' CR/CAR narrows the admissible family of constitutions but selects no substantive content.
\end{itemize}
\end{paperbox}

\begin{paperbox}
\textbf{\cb{} --- Systematic Restatements.} A result is in \cb{} if the consequence is contained in its premise, so that the ``derivation'' is a restatement of what the premise already says. \cb{} results are not framework contributions; they are labels that make the framework's content explicit.
\begin{itemize}
\item \textbf{F2 (Agency-Reciprocity) from B4.} B4 says agency is a complete ground for standing. F2 says agents have equal standing. F2 is contained in B4.
\item \textbf{F3 (Defensive Content) from B5.} B5 says standing constrains subtraction and does not license predation. F3 says the same thing.
\item \textbf{T1 (Ground-Indexed Consistency) from L1.} T1 applies L1 to the case of a claimant that holds a claim on a ground for itself.
\item \textbf{T2-D (Discursive Non-Interference) from T1 and B6--B7.} T2-D restates the protection of the agency preconditions in the terms already used by B6.
\item \textbf{T3 (Epistemic Integrity) from T2-D and B6(b).} T3 restates the protection of accurate information.
\item \textbf{T4 (Defense Against Initiatory Acts) from T2-D, T2-G, and INIT.} T4 restates the consequence of the definition of an initiatory act.
\item \textbf{T5 (Rejection of Retribution) from T4, L1, and B8.} T5 restates the consequence of harm to an N-conduct agent together with B8.
\item \textbf{T7a (Property Consistency), CE (Complaint Estoppel), CSL (Commitment-Scope Lemma).} Each applies L1 to a specific structural case.
\end{itemize}
\end{paperbox}

\textbf{What this means for the framework's contribution.} The framework's object-level derivational content is exactly two results, L1 and NRD, together with the meta-level constraint CR/CAR. Its contribution beyond those results is not derivational; it is the \emph{audit architecture}---the discipline of naming every premise, pricing every rejection, publishing the dependency graph, and disclosing the residual basis. The framework's value is therefore (i) the non-trivial derivations, (ii) the audit architecture, and (iii) the runtime enforcement design. It is not the derivation of the ethics from first principles, and it should not be read as claiming that.

A reader who is looking for a nontrivial derivation of F2 or F3 from thinner premises will not find one, because none exists. The reader who is looking for a specific, auditable, and honestly bounded localization of a robust practice will find what they need.

\subsection{The Three-Level Practice Structure}
The framework's practice is characterized at three levels, with a fourth reflexive refinement.
\begin{longtable}{@{}L{1.3cm} L{4.2cm} L{9.4cm}@{}}
\toprule
\textbf{Level} & \textbf{Name} & \textbf{Constitutive Rules} \\
\midrule
\endfirsthead
\toprule
\textbf{Level} & \textbf{Name} & \textbf{Constitutive Rules} \\
\midrule
\endhead
\bottomrule
\endfoot
\textbf{1} & Thin reason-giving & Offer, Object, Respond. Minimal speech act of presenting a ground, under ARC. \\
\textbf{2} & Accountable reason-giving & Level 1 plus RR. Yields universalization of stated grounds (L1). \\
\textbf{3} & Robust reciprocal reason-giving & Level 2 plus AR and the basic constitutive standards (B4--B8). Yields the core ethics. \\
\textbf{3R} & Reflexively admissible robust practice & Level 3 plus CR/CAR: constitutional rules are themselves answerable; self-authenticating authority is excluded, while substantive plurality remains. \\
\end{longtable}

\textbf{What Level 3 and Level 3R are and are not.} Level 3 is a specific robust practice. It is one coherent robust practice among several. The thin practice does not select Level 3 uniquely; multiple robust practices are coherent, differing in their basic constitutive standards. Level 3R is the reflexive admissibility layer added to Level 3: it narrows the admissible family by prohibiting a constitution from using its own challenged authority as its sole justification. The framework adopts Level 3 as its specification. It does not claim that Level 3 is the unique extension of Level 2, and it does not claim that agents who participate in Level 2 are thereby committed to Level 3. The framework's claim is precisely this: agents that participate in \emph{this specified} Level~3 practice are committed to the ethics derived from it.

\subsection{PRAC as an Idealization}
B7 (PRAC) requires a note. It states that the practice of giving and seeking reasons presupposes participants who can deliberate and receive reasons uncoerced and accurately informed. This is an idealization of an ideal practice, not a description of empirical practices of reason-giving.

\textbf{Why the idealization.} Real reason-giving practices tolerate coercion, deception, and manipulation. Participants in real discourse are often non-ideal: they are subject to incentives, incomplete information, and social pressure. Liars participate in assertion without satisfying PRAC. Under the framework's specification, B7 describes the practice the framework analyzes---the ideal practice of addressed justification in which participants can genuinely deliberate and receive reasons. The framework does not claim that real practices satisfy B7. It claims that the ethics it derives is the ethics of the ideal practice, and that real practices approximate it to varying degrees.

\textbf{Consequence for the framework's scope.} Real deployments approximate the ideal practice and are covered to the extent that they do. An operator that deploys a system in a context where coercion or deception is the norm has a correspondingly weaker coverage by the framework's core ethics. The framework states this openly.

\subsection{Practice Selection}\label{sec:selection}
The framework is one coherent robust practice among several. A hierarchical practice with role-sensitive standing, an enforcement-bearing practice, or a practice that treats agency as one of several grounds for standing are each coherent alternatives. The framework does not claim to select among them by logic. It therefore faces an honest question: why this practice rather than another?

\textbf{The hidden premise made explicit.} The premise that does the selecting is: \emph{this practice ought to be the one deployed.} The framework does not resolve this premise. It cannot resolve it by derivation. The premise is a decision.

\begin{paperbox}
\textbf{The Practice-Selection Decision.} The framework treats the practice-selection decision as belonging to the deploying operator. It is made public through the operator's claims to its users, its regulators, and the public. The framework's role is to check that the deployed practice is consistent with those claims, not to adjudicate which practice is correct.
\end{paperbox}

\textbf{Why the operator's public claims are the right locus.} The framework's ethics is constitutively grounded in the practice of making claims. When an operator publicly claims that its system treats all users fairly, refuses harmful requests, or respects user autonomy, the operator has entered the practice. The operator's own claims commit it to a specific robust practice: the one whose basic standards the operator's claims presuppose. The framework's role is to check the operator's deployments against the operator's claims.

\textbf{Practice selection under the Claims Registry.} Practice selection remains a decision of the deploying operator, but the selected practice must be entered into the Claims Registry (Section~\ref{sec:governance}), bound to the active \SafetyBundle{}, and accompanied by evidence sufficient to support the assurance label claimed for the deployment.
\begin{enumerate}
\item \textbf{The selected practice is explicit.} Public claims are mapped to the substantive standards they invoke; ambiguity is recorded rather than silently resolved in favor of stronger assurance.
\item \textbf{The selected practice is versioned.} Changes are cryptographically bound to a new constitution, profile, or policy version according to their scope.
\item \textbf{Material changes are governed.} A change that relaxes or enlarges the safety floor follows the non-monotone update path in Section~\ref{sec:governance}; restrictive updates may follow the monotone path.
\item \textbf{The evidence scope controls the claim.} If the selected practice cannot be checked over a material deployment surface, the deployment label is downgraded to the weakest supported assurance class.
\end{enumerate}

\textbf{What the framework does not do.} The framework does not decide that equal treatment is the correct practice. It does not adjudicate between the framework's practice and its coherent alternatives. It does not hold operators to a specific ethical commitment. It checks that the operators' commitments, whatever they are, are internally consistent across their public claims and their deployed systems.

\textbf{The limits of the claims-to-practice mapping.} The framework does not publish a full mapping from possible public claims to selected practices. Some claims---for example, ``our system is helpful''---are too vague to select a specific practice. Some claims---for example, ``our system treats all users fairly''---are specific enough to commit the operator to a practice with B4 and B5. Where the operator's claims are vague, the framework cannot check consistency, and the operator remains free to select among practices that would satisfy the vague claims. This is a limitation of the framework's check, not a claim that vague claims are ethically acceptable. A separate legal or regulatory framework would be required to constrain operators whose public claims are deliberately vague.

\textbf{The family of practices.} Because the framework does not select a unique practice, it should be understood as one explicitly specified member of a larger family of robust practices. The present specification is the egalitarian, defensive, agency-complete practice. Institutional-authority and role-sensitive practices can be represented as alternative constitutions by changing the declared substantive standards while retaining the same reflexive admissibility and runtime machinery.

\textbf{The audit trail for practice selection.} The dependency audit (Section~\ref{sec:rab}) covers inference; the Claims Registry and its checkability statuses cover decision. Both are required for full auditability. The registry records what was selected and why, but it cannot certify that the selection is correct; that judgment remains with the deploying operator and the external reviewers named in Section~\ref{sec:governance}.

\subsection{Constitutional Pluralism Audit}
The framework should be tested not only against its own constitution but against alternatives. The point is not to prejudge the result; it is to identify exactly which disagreement remains after reflexive closure.

\begin{longtable}{@{}L{3.1cm} L{3.5cm} L{3.8cm} L{4.4cm}@{}}
\toprule
\textbf{Candidate constitution} & \textbf{Can satisfy thin practice?} & \textbf{Can satisfy \CR{}/\CAR{}?} & \textbf{Where disagreement remains} \\
\midrule
\endfirsthead
\toprule
\textbf{Candidate constitution} & \textbf{Can satisfy thin practice?} & \textbf{Can satisfy \CR{}/\CAR{}?} & \textbf{Where disagreement remains} \\
\midrule
\endhead
\bottomrule
\endfoot
Agency-complete, defensive (this framework) & Yes & Yes & Chooses B4--B8. \\
Role-sensitive institutional constitution & Yes & Potentially yes & Role/office may remain a primitive substantive standard; B4 is rejected. \\
Enforcement-bearing constitution & Yes & Potentially yes & Standing may include authorized coercive powers; B5 is rejected. \\
\end{longtable}

This audit is strategically important. It prevents the strengthened result from being overstated as a uniqueness theorem. The contribution is a narrowing theorem: reflexive reason-giving excludes self-authenticating constitutional authority, but it does not by itself collapse all substantive constitutional plurality.

\subsection{Formal Ground Calculus and Operational ARC}\label{sec:ground}
The reason-giving vocabulary needs a semantic object that is stable before an objection is evaluated. Otherwise an objection can retroactively redefine the alleged ground until the claim appears either harmless or self-justifying. The paper therefore introduces a typed \Ground{} object.
\begin{paperbox}
\textbf{Ground object.} A ground is a typed tuple
\[
g=\langle id,\; predicate,\; domain,\; binding,\; evidence,\; scope,\; validity,\; effect,\; provenance\rangle.
\]
The \emph{predicate} states the condition appealed to; \emph{domain} states whose treatment the ground purports to distinguish; \emph{binding} records the relevant role/principal/resource parameters; \emph{evidence} identifies what supports the predicate; \emph{scope} specifies where the ground applies; \emph{validity} specifies temporal and evidentiary validity; \emph{effect} identifies the normative or operational consequence; and \emph{provenance} identifies who introduced the ground and from which policy or constitutional artifact.
\end{paperbox}

\textbf{No-Retrofitting Rule (RR-G).} Once a challenge is raised, the claimant may clarify a ground only by operations that preserve its recorded predicate, domain, binding, scope, and intended effect. A material change creates a new ground and cannot be used to discharge the earlier challenge. This rule is required to make L1 testable rather than dependent on post-hoc semantic repair.

For ARC, the phrase ``could accept'' is operationalized in three separable predicates. A ground is \emph{Addressable} when its predicate and scope can be stated to the addressee in a role-independent form. It is \emph{Evaluable} when the addressee or an authorized evaluator can inspect the evidence relevant to the predicate. It is \emph{Defensible} when the ground survives a predefined challenge set including publicity, idealized-counterpart, role-reversal, and adversarial-panel tests. The framework does not assert that any empirical panel proves moral truth; the tests establish whether the declared practice applies its own acceptance standard consistently.

L1 therefore has two components. \emph{L1-Public} is the invariant claim that a differentiating reason must be an addressable public ground. \emph{L1-Relevance} is stronger: the ground must actually bear on the normative difference. The \ADM{} analysis (Section~\ref{sec:adm0}) establishes the first component as an invariant of the specified thin practice; relevance is carried by the declared answerability and robustness premises. This split prevents the paper from presenting relevance as a consequence of admissibility alone.

\subsection{Constitutional Safety as a Typed Interface}
A constitution is not itself a runtime command. To become executable it must instantiate the typed safety ontology defined in Section~\ref{sec:engine}. The interface is therefore:
\[
\text{Constitution}\rightarrow\text{Safety Profile}\rightarrow\text{Policy}\rightarrow\text{Proof}\rightarrow\text{Capability}\rightarrow\text{Effect}.
\]
At each arrow the type of object changes. An argument about normative authority is not silently treated as a predicate about machine state; a machine predicate is not silently treated as empirical evidence; and a signed capability is not treated as proof that the underlying constitution is ethically correct. This typed separation is the central epistemic firewall of the paper.

\section{The Residual Axiom Basis}\label{sec:rab}
\subsection{The Basis, Stated Precisely}
The framework's foundation is the set of basic constitutive standards named in Section~\ref{sec:robust}. They are restated here with their tier, their status, and their function.
\begin{longtable}{@{}L{1.0cm} L{4.9cm} L{0.9cm} L{1.4cm} L{7.0cm}@{}}
\toprule
\textbf{ID} & \textbf{Statement} & \textbf{Tier} & \textbf{Status} & \textbf{Function} \\
\midrule
\endfirsthead
\toprule
\textbf{ID} & \textbf{Statement} & \textbf{Tier} & \textbf{Status} & \textbf{Function} \\
\midrule
\endhead
\bottomrule
\endfoot
B1 & \Atwo{}: a claim offers a justificatory reason to the addressee. & 1a & \statusN{}\textsuperscript{*} & Constitutive of claiming. \\
B1$'$ & ARC: a ground is a ground only if it is a reason the addressee could accept; private causes are not grounds. & 1a & \statusN{}\textsuperscript{*} & Constitutive of addressed justification. Defeats the agent-relative countermodel to L1; load-bearing for F2. \\
B2 & RR: a challenged justification must be defended, qualified, or withdrawn. & 1a & \statusN{}\textsuperscript{*} & Constitutive of accountable reason-giving. \\
B3 & AR: participants are answerable for the normative distinctions they introduce. & 1b & \statusN{} & Constitutive of robust reciprocal reason-giving, or substantive premise. \\
B4 & Agency-Completeness: agency is a complete ground for basic standing. & 1b & \statusN{} & Terminates the regress. Enables F2. \\
B5 & Defensive Interpretation: standing constrains subtraction; it does not license predation. & 1b & \statusN{} & Enables F3. \\
B6 & A3 and Basic Goods: (a) freedom from coercive override, (b) access to accurate information, (c) capacity to deliberate, (d) Basic goods (life, bodily integrity). & 2 & \statusN{} & Declared list; contestable. \\
B7 & PRAC: the practice presupposes participants who can deliberate and receive reasons uncoerced and accurately informed. & 2 & \statusN{} & Load-bearing for T2-D and T3. \\
B8 & Forward-Looking Grounds: reasons for inflicting harm on an agent must have a forward-looking component. Reasons that refer only to an agent's past conduct are idle. & 2 & \statusN{} & Enables T5 (rejection of retribution). \\
\end{longtable}
\textbf{Tier 1a} items are constitutive of the practice at levels 1 and 2. Their denial is coherent only by not participating in the respective practice. \textbf{Tier 1b} items are the basic standards that terminate the level-3 regress. Their denial yields a different robust practice. \textbf{Tier 2} items are declared and contestable. B1$'$ is the companion of B1 and is counted with it in the eight-item basis below.

\subsection{Necessary versus Declared Status}
The framework distinguishes two kinds of \statusN{} status.
\begin{paperbox}
\textbf{Status \statusN{}\textsuperscript{*} (Necessary).} A premise is \statusN{}\textsuperscript{*}-necessary if denying it while continuing to participate in the practice is incoherent. B1 (\Atwo{} + ARC) and B2 (RR) are the necessary premises. Their denial is coherent only by not claiming or not being accountable. They are the single points of failure of the framework: dropping either removes every core result.

\textbf{Status \statusN{} (Declared).} A premise is \statusN{}-declared if it is a basic standard of the robust practice that the framework adopts, and its denial is coherent but yields a different robust practice. B3--B8 are declared. Their denial changes the practice; it does not exit the practice. B3 (AR) is on the boundary: a participant who rejects AR has moved from the robust practice back to the thin practice, which is a coherent alternative but not the framework's.
\end{paperbox}

\textbf{Consequence for the framework's central claim.} The central claim rests on B1 and B2, which are necessary, and on B3--B8, which are declared. The framework does not claim that B3--B8 are unavoidable. It claims that they are its chosen standards, stated explicitly so that a reader can see exactly what is being assumed and what follows.

\textbf{A third distinction: invariance.} Beyond the necessary/declared split, the framework distinguishes which of its claims are \emph{invariant} across admissible reason-giving practices and which are \emph{constitution-specific}. Subsection~\ref{sec:adm0} establishes that under a minimal admissibility semantics, only the public-difference component of L1 is invariant; B3--B8 are constitution-specific. This is the framework's central formal finding, and it is what distinguishes the framework's ethics from a claim of universal morality.

\subsection{The Invariant Floor Search}\label{sec:adm0}
The framework's basic constitutive standards B1--B8 are declared. The question the audit asks is not whether they are declared (they are) but whether they are \emph{invariant} across the class of admissible reason-giving practices: whether any practice that satisfies the constitutive requirements of public, auditable reason-giving entails them. A finite countermodel search answers this question. The answer is: no, not under a minimal admissibility semantics.

\subsubsection{The Minimal Admissibility Semantics (\ADM{})}
\begin{paperbox}
\textbf{\ADM{} (Minimal Admissibility).} A finite constitution $C$ is \ADM{}-admissible when its normative outputs satisfy four constraints:
\begin{enumerate}[label=(\alph*)]
\item \textbf{Publicity.} Normative outputs depend only on public ground/role facts.
\item \textbf{No Hidden Identity.} Private claimant identity and private motive are not admissible normative inputs.
\item \textbf{Public Grounding.} Every normative difference is indexed to an explicit public differentiator.
\item \textbf{No Substantive Axiom.} \ADM{} does not contain B4--B8 or any other substantive ethical premise.
\end{enumerate}
\end{paperbox}
\ADM{} formalizes public, auditable justification. It is the floor that any practice of giving and asking for reasons with public accountability is expected to satisfy. It is not a complete formalization of natural-language reason-giving; it is a deliberately minimal semantics chosen so that the search does not simply rediscover the framework's own axioms.

\subsubsection{Finite Domain and Search Method}
The search uses a finite domain: two agents ($a$, $b$), two public roles (\texttt{ordinary}, \texttt{authority}), two public ground classes ($G$, $H$). A normative rule is a function from public states to treatment values. For each candidate proposition $P$, the search attempts to construct a finite \ADM{}-admissible constitution in which $P$ fails. Exhaustive enumeration is possible for the small domain; the L1 search checks $256$ public-state/treatment configurations.

\subsubsection{Countermodel Results}
\begin{longtable}{@{}L{3.9cm} L{5.6cm} L{6.4cm}@{}}
\toprule
\textbf{Proposition} & \textbf{\ADM{} Result} & \textbf{Status} \\
\midrule
\endfirsthead
\toprule
\textbf{Proposition} & \textbf{\ADM{} Result} & \textbf{Status} \\
\midrule
\endhead
\bottomrule
\endfoot
L1 (public-difference component) & No countermodel in 256 checked configurations & Invariant \\
L1 (relevance component) & Countermodel & Not invariant \\
NRD (full form) & Countermodel & Not invariant \\
B4 (Agency-Completeness) & Countermodel & Constitution-specific \\
B5 (Defensive Interpretation) & Countermodel & Constitution-specific \\
B6 (A3 and Basic Goods) & Countermodel & Constitution-specific \\
B7 (PRAC) & Countermodel & Constitution-specific \\
B8 (Forward-Looking Grounds) & Countermodel & Constitution-specific \\
T2-D (Discursive Non-Interference) & Countermodel & Constitution-specific \\
T3 (Epistemic Integrity) & Countermodel & Constitution-specific \\
T4 (Defense Against Initiatory Acts) & Countermodel & Constitution-specific \\
T5 (Rejection of Retribution) & Countermodel & Constitution-specific \\
\end{longtable}
Every countermodel is a \emph{publicly specified, role-indexed rule}. None exploits hidden identity. None relies on private motive. Each is a coherent public constitution that \ADM{} admits.

\textbf{The structural defect shared by the countermodels.} Each exploits a role asymmetry: the same public ground receives different normative force depending on the role of the claimant or addressee. \ADM{} does not prohibit this. It requires only that the role difference be public.

\subsubsection{The Two-Tier Architecture}
\begin{paperbox}
\textbf{The Invariant Floor.} Under \ADM{}, the invariant floor contains:
\begin{itemize}
\item \textbf{Public-ground consistency.} L1's public-difference component: same public ground, same treatment, unless a public differentiator is supplied.
\item \textbf{Public-difference accountability.} Any normative difference must be indexed to an explicit public differentiator.
\end{itemize}
The invariant floor is thin but formal. Its strongest demonstrated theorem is L1's public-difference component.

\medskip
\textbf{The Declared Ethics.} B4--B8 and their consequences (T2-D--T6, T7a, T8, T11a, T16, CE, CSL) are declared standards, not invariant under \ADM{}. They are not made invariant by Normative Role Closure (NRC): the NRC search (below) shows that the residual countermodels are symmetric constitutions, immune to any constraint stated in role vocabulary.
\end{paperbox}

\textbf{The framework's claim is not that the ethics is invariant.} The framework's claim is that the ethics is:
\begin{itemize}
\item Derived from the declared standards (\ca{}) or restated from them (\cb{}).
\item Audited against the standards (dependency audit, Appendix~\ref{app:audit}).
\item Not invariant under the minimal admissibility semantics.
\item Enforced architecturally where it applies (Section~\ref{sec:engine}).
\end{itemize}
The distinction between invariant and declared is the framework's central formal finding. It is what the framework delivers that a claim of universal morality would not.

\subsubsection{The NRC Search and Its Result}\label{sec:nrc}
The next formal search tested whether \textbf{Normative Role Closure} (NRC) eliminates the countermodels to B4--B8 without itself containing B4--B8. The search has been run over the finite domain above. This subsubsection reports the result.

\begin{paperbox}
\textbf{Normative Role Closure (NRC).} A public institutional role may alter \emph{who} performs an act, but may not by itself alter the basic normative force of a ground, unless the role difference is itself justified by a ground that remains valid under role reversal.
\end{paperbox}

\textbf{Formalization.} In the two-role domain, NRC admits two readings. Under the \emph{weak} reading (role-swap symmetry), treatment is invariant under exchange of claimant and addressee roles: $C(\mathrm{ord},\mathrm{auth},g) = C(\mathrm{auth},\mathrm{ord},g)$. Under the \emph{strong} reading (role independence), the normative force of a ground may not depend on roles at all: since roles are the only non-ground public facts in the domain, no role difference can be justified by a role-reversal-valid ground, so $C$ must be a function of the ground class alone. The strong reading subsumes the weak; the result holds under both.

\textbf{Method: lemma first, enumeration second.} In this domain the NRC result does not rest on enumeration. Under the strong reading, every admissible constitution is a function of the ground class alone, and there are exactly four such functions: $\{G, H\} \to \{0, 1\}$. The result is established analytically (Observation below). The 256-configuration enumeration above was re-run with NRC added as a fifth admissibility constraint and agrees with the lemma; the enumeration is a consistency check on the encoding, not the basis of the claim.

\begin{longtable}{@{}L{4.6cm} L{2.5cm} L{2.6cm} L{4.9cm}@{}}
\toprule
\textbf{Proposition} & \textbf{\ADM{}} & \textbf{\ADM{}+NRC} & \textbf{Status} \\
\midrule
\endfirsthead
\toprule
\textbf{Proposition} & \textbf{\ADM{}} & \textbf{\ADM{}+NRC} & \textbf{Status} \\
\midrule
\endhead
\bottomrule
\endfoot
L1 (public-difference component) & 0 countermodels & 0 countermodels & Invariant \\
B4 (Agency-Completeness) & 240 countermodels & 2 countermodels & Constitution-specific \\
B5 (Defensive Interpretation) & 240 & 2 & Constitution-specific \\
B6 (A3 and Basic Goods) & 240 & 2 & Constitution-specific \\
B7 (PRAC) & 240 & 2 & Constitution-specific \\
B8 (Forward-Looking Grounds) & 240 & 2 & Constitution-specific \\
T2-D (Discursive Non-Interference) & 240 & 2 & Constitution-specific \\
T3 (Epistemic Integrity) & 240 & 2 & Constitution-specific \\
T4 (Defense Against Initiatory Acts) & 240 & 2 & Constitution-specific \\
T5 (Rejection of Retribution) & 240 & 2 & Constitution-specific \\
\end{longtable}

Under the strong reading, exactly four constitutions remain admissible, and every one of B4--B8 and T2-D--T5 is falsified by at least two of them. NRC therefore does not eliminate the countermodels.

\begin{paperbox}
\textbf{Observation (Symmetric Countermodel Lemma).} Let $\mathcal{S}$ be any admissibility semantics whose constraints are invariant under permutations of public roles. Then no proposition that constrains the treatment of a ground class---as B4--B8 and T2-D--T5 each do---is invariant under $\mathcal{S}$.

\emph{Proof.} Role-invariant constraints cannot distinguish constitutions that are already uniform across roles. The four NRC-strong admissible constitutions are pairwise distinguished only by the treatment they assign to ground classes. A constraint that excludes the nihilist constitution (no standing on any ground) must assert that some ground succeeds for someone; a constraint that excludes the uniformly predatory constitution (every agent's bare standing licenses override of every other) must assert that some ground does not license override. Either assertion is substantive: it specifies what a ground normatively does, which is exactly the content of B4--B8. Since $\mathcal{S}$, by hypothesis, carries no substantive content, no such exclusion is available, and each proposition is falsified by at least one $\mathcal{S}$-admissible constitution.\hfill$\square$
\end{paperbox}

\textbf{The third outcome.} The search was pre-registered as a dichotomy: either NRC thickens the floor, or it succeeds only by restating B4--B8. The result is a third outcome. NRC is \emph{independent} of B4--B8: of the four NRC-admissible constitutions, exactly half satisfy each proposition, so NRC neither entails nor restates them. And NRC is \emph{insufficient}: the substantive ethics remains falsifiable within its admissible class. NRC eliminates exactly the role-privilege countermodels and nothing else.

\textbf{Consequences for the invariant floor.} Adding NRC to \ADM{} yields a third invariant alongside public-ground consistency and public-difference accountability:
\begin{paperbox}
\textbf{Role-Blindness (invariant under \ADM{}+NRC).} No constitution is admissible in which the normative force of a ground depends on the public roles of claimant or addressee. Bare role privilege---the same public ground carrying different normative force by role alone---is inadmissible.
\end{paperbox}
The floor is thicker. It does not reach the ethics. The residual countermodels are symmetric constitutions---egalitarian in form, non-egalitarian in content---and no constraint stated in role and publicity vocabulary can exclude them, because they violate nothing that vocabulary can express. Excluding them requires substantive content: exactly B4--B8. The two-tier architecture is therefore structural, not provisional.

\textbf{Status of the search:} run and reported. The strong reading of NRC is adopted as the specification. The theorem is the Symmetric Countermodel Lemma, established analytically; the countermodel counts are encoding artifacts of the two-role, two-ground domain and carry no independent weight. Enlarging the domain increases the counts and cannot change the lemma, whose proof does not depend on domain size. The encoding used for the consistency-check enumeration is declared, not audited; the lemma is what the framework rests on.

\subsection{Intermediate Results}
Before the core ethics can be derived, four intermediate results must be established.

\subsubsection{ARC as a Constitutive Companion of \Atwo{}}
ARC is stated as B1$'$, the companion of \Atwo{} (B1). \Atwo{} says a claim offers a justificatory reason to the addressee. ARC specifies what counts as a justificatory reason: a ground the addressee could accept as bearing on the claim. Without ARC, \Atwo{} can be read weakly, and the weak reading reduces \Atwo{} to a description of the claimant's psychology. The framework adopts ARC as part of what \Atwo{} means. \textbf{Status of ARC:} constitutive premise, listed as B1$'$ in the table above. Results that depend on it are so labeled in the dependency audit.

\subsubsection{\Ptwo{}: The Ground Requirement}
\begin{paperbox}
\textbf{\Ptwo{} (Ground Requirement).} Differential treatment of agents requires a ground that can be stated and defended to those treated.
\end{paperbox}
\textbf{Derivation.} Immediate from B1. A claim invokes a standard offered to the addressee. If the claim differentiates between agents, the differentiating standard must be offered. The standard is the ground. A ground that cannot be stated and defended to the addressee is not a claim's ground; it is at most a private cause.
\textbf{Status:} derived (\statusD{}) from B1.

\subsubsection{L1: The Universality of Grounds}
L1 is stated in Section~\ref{sec:accountable}. Its derivation from B1 + ARC + B2 is: B1 with ARC makes the ground addressable to the addressee; B2 requires that a challenged ground be defended, qualified, or withdrawn; combining the two yields the universalization proposition, of which L1 is the general form. \textbf{Status:} derived (\statusD{}), \ca{}.

L1 has two components:
\begin{enumerate}[label=(\roman*)]
\item \textbf{Public-difference accountability.} Any differentiating reason must be publicly stateable. This is the requirement that the reason be an addressable ground in the sense of ARC.
\item \textbf{Relevance.} The differentiating reason must actually bear on the appropriateness of the claim. This is the requirement that the reason not be idle.
\end{enumerate}
Under the \ADM{} search, \emph{only component (i) is invariant}. Component (ii) requires an additional premise, and it is where NRD declares its relevance requirement. The framework therefore does not claim that the full L1 is invariant under \ADM{}. It claims that L1's public-difference component is invariant and that its relevance component is declared.

\subsubsection{NRD: No Unjustified Normative Difference}
NRD is stated in Section~\ref{sec:robust}. Its derivation is from B3 (AR) and B2 (RR): B3 makes every normative difference subject to the practice's justificatory demand; B2 specifies what the justificatory demand is. \textbf{Status:} derived (\statusD{}), \ca{}.

\subsection{The Full Dependency Audit}\label{sec:rab:audit}
Every core result of the framework rests on a set of basic constitutive standards computed by tracing the declared dependency graph to its leaves. The audit includes the derived definitions (INIT, ONG, N-conduct, stakes rank) and the declared principles (FSP, RW, O1) used by the core results.
\begin{longtable}{@{}L{2.6cm} L{7.4cm} L{5.3cm}@{}}
\toprule
\textbf{Result} & \textbf{Rests on (non-derived premises)} & \textbf{Survives rejection of B4, B5?} \\
\midrule
\endfirsthead
\toprule
\textbf{Result} & \textbf{Rests on (non-derived premises)} & \textbf{Survives rejection of B4, B5?} \\
\midrule
\endhead
\bottomrule
\endfoot
\Ptwo{} & B1 & Yes \\
L1 & B1, B2 & Yes \\
NRD & B3, B2 & Yes \\
T1 (Ground-Indexed Consistency) & B1, B2 & Yes \\
CE, T7a, CSL & B1, B2 & Yes \\
T2-D, T3 & B1, B2, B6(a,b,c), B7 & Yes \\
F2 & B1, B2, B3, B4 & No \\
T2-G & B1, B2, B3, B4, B5 & No \\
F3 & B1, B2, B3, B4, B5 & No \\
T4 & B1, B2, B3, B4, B5, B6, B7 & No \\
T5 & B1, B2, B3, B4, B5, B6, B7, B8 & No \\
T6 & B1, B2, B6, O1 & Yes \\
T8 (Partiality) & FSP & Yes \\
T11a (Return of Goods) & B1, B2 & Yes \\
T16 (No Group Liability) & B1, B2, RW & Yes \\
\end{longtable}
The \textbf{Residual Axiom Basis} for the full core ethics (T1--T8, T11a, T16, CE, CSL) is therefore $\{$B1, B2, B3, B4, B5, B6, B7, B8$\}$, with P1 (Identity) in the background. Four findings follow.
\begin{enumerate}
\item \textbf{The basis is eight items.} Every core result is either derived from these eight or is one of them. Each is stated, argued, and priced.
\item \textbf{B1 and B2 are the single points of failure.} Dropping B1 or B2 removes every core result.
\item \textbf{The declared standards carry the specific ethics.} B4, B5, and B8 are what make the framework an egalitarian, defensive, non-retributive practice. Their denial yields a different robust practice.
\item \textbf{The most contested standards carry the narrowest results.} Ground-indexed consistency (T1), complaint estoppel (CE), property consistency (T7a), and the commitment-scope lemma (CSL) survive the rejection of B4, B5, and B8. A participant that rejects the framework's egalitarian or defensive standards remains bound by these results. B4, B5, and B8 matter only for general agency non-interference, defense, no-retribution, and necessity.
\end{enumerate}
\textbf{Limits of the audit.} The dependency graph is a check of the declared structure. It shows what the stated derivations use. It does not show that the derivations are valid; that is a matter for machine-checking (Section~\ref{sec:status}).

\newpage
\subsection{The Decomposition Chain}
\begin{paperbox}
\textbf{The Decomposition Chain.}
\[
\begin{aligned}
&A=A &&\text{(Identity; no ethics follows)}\\
&\Downarrow\\[-0.2em]
&\Atwo{}+\mathrm{ARC}+\mathrm{RR} &&\text{(addressed, responsive reason-giving)}\\
&\Downarrow\\[-0.2em]
&\mathrm{L1}_{\mathrm{public}}+\mathrm{T1}+\mathrm{CE}+\mathrm{T7a}+\mathrm{CSL}\\
&\Downarrow\\[-0.2em]
&\mathrm{AR}+\mathrm{NRD}+\mathrm{CR/CAR} &&\text{(answerability for normative difference)}\\
&\Downarrow\\[-0.2em]
&\mathrm{B4\text{--}B8} &&\text{(declared substantive standards)}\\
&\Downarrow\\[-0.2em]
&\text{Robust reciprocal practice}\rightarrow\text{Core ethics}
\end{aligned}
\]
Every arrow is either a derivation or an explicitly declared addition; no arrow silently changes the status of its premises.
\end{paperbox}
Every arrow in the chain is either a derivation (\statusD{}) or a stated constitutive addition. There are no hidden steps. The framework is explicit that B3--B8 are the points at which the chain stops deriving and starts declaring.

\section{The Core Ethics}\label{sec:core}
\subsection{What the Core Ethics Is}
Before listing the results, the framework states the substantive content of the ethics plainly, to prevent a reader from mistaking the machinery of L1 and NRD for the ethics itself.
\begin{paperbox}
\textbf{The Framework's Ethics, Stated Substantively.} The framework's ethics is a \emph{negative-rights, non-aggression core}. It prohibits:
\begin{itemize}
\item \textbf{Initiatory interference} with another agent's basic goods (life, bodily integrity) or agency preconditions (freedom from coercive override, access to accurate information, capacity to deliberate).
\item \textbf{Deception} that attacks another agent's access to accurate information.
\item \textbf{Obstruction of deliberation} that attacks another agent's capacity to deliberate.
\item \textbf{Unjustified normative differentiation} that is not supported by a relevant ground.
\item \textbf{Retribution}, meaning harm inflicted on an agent only because of past harm.
\end{itemize}
It does not deliver:
\begin{itemize}
\item \textbf{Positive duties}: it does not require agents to aid, rescue, or assist others. A general duty to aid is not derivable from the framework's core.
\item \textbf{Distributive justice}: it does not deliver a principle of distribution of resources or opportunities.
\item \textbf{Harm to non-agents}: animals, ecosystems, and beings below the standing threshold are outside the base framework's scope.
\item \textbf{Positive rights}: it does not deliver entitlements to goods or services.
\end{itemize}
The framework's ethics is therefore a specific political morality: a negative-rights, non-aggression core with no positive duties, no distributive principles, and no harm-to-non-agents coverage. This is stated openly rather than left implicit.
\end{paperbox}
The framework's ethics is one coherent robust practice among several. A virtue-ethical practice, a consequentialist practice, and a positive-rights practice are all coherent alternatives. The framework's contribution is not to adjudicate among them; it is to specify the negative-rights, non-aggression core with an audit architecture that makes its premises and its limits visible.

\subsection{The Core Results Table}
Each result is stated with its status (\ca{} or \cb{}) and its premise dependencies.
\begin{longtable}{@{}L{1.2cm} L{5.6cm} L{1.9cm} L{6.4cm}@{}}
\toprule
\textbf{ID} & \textbf{Result} & \textbf{Status} & \textbf{Depends On} \\
\midrule
\endfirsthead
\toprule
\textbf{ID} & \textbf{Result} & \textbf{Status} & \textbf{Depends On} \\
\midrule
\endhead
\bottomrule
\endfoot
T1 & Ground-Indexed Consistency & \statusD{}, \cb{} & B1 + B2 (via L1) \\
T2-D & Discursive Non-Interference & \statusD{}, \cb{} & T1, B6(a,b,c), B7 \\
T2-G & General Agency Non-Interference & \statusD{}, \cb{} & T1, F2, F3 \\
T3 & Epistemic Integrity & \statusD{}, \cb{} & T2-D, B6(b) \\
T4 & Defense Against Initiatory Acts & \statusD{}, \cb{} & T2-D, T2-G, INIT (definition) \\
T5 & Rejection of Retribution & \statusD{}, \cb{} & T4, L1, B8 \\
T6 & Necessity & \statusD{}, \cb{} & T1, B6, O1 \\
T7a & Property Consistency & \statusD{}, \cb{} & T1 \\
T8 & Partiality & \statusD{}, \cb{} & FSP \\
T11a & Return of Goods Still Held & \statusD{}, \cb{} & T1, ONG (definition) \\
T16 & No Group Liability & \statusD{}, \cb{} & CE, RW \\
CE & Complaint Estoppel & \statusD{}, \cb{} & B1, L1 \\
CSL & Commitment-Scope Lemma & \statusD{}, \cb{} & B1, L1 \\
\end{longtable}
The core results are all \cb{}; the derivational work is done by L1 and NRD. This is expected: the framework's contribution is localized in L1, NRD, and CR/CAR, not in the specific rules that follow from the basic standards.

\subsection{Foundational Definitions}
Three definitions and one stakes ranking are used by the core results.
\subsubsection{A3 Impact, N-Conduct, and Voluntary Conduct}
\begin{paperbox}
\textbf{A3 Impact.} A conduct has an \emph{A3 impact} on an agent $Y$ if and only if it coercively overrides or destroys $Y$'s execution of a chosen action, deceives $Y$ about a standard in play, or incapacitates $Y$'s deliberation. The definition is descriptive: it does not invoke ``protection'' or ``violation.'' It identifies conduct by its causal effect on another agent.

\textbf{Voluntary Conduct.} An agent's conduct is \emph{voluntary} if it is performed on a ground. Conduct on autopilot, in a reflex, or under coercion is not voluntary in this sense.

\textbf{N-Conduct.} An agent's conduct is \emph{N-conduct} if it has no A3 impact on any other agent. It includes speaking, moving, holding a resource, and refusing. It does not include conduct that overrides, deceives, or incapacitates.
\end{paperbox}
These definitions provide the terms INIT, T4, T5, and O0 use. They are descriptive and carry no moral content themselves.

\subsubsection{INIT: The Initiation Test}
\begin{paperbox}
\textbf{INIT (Initiation Test).} An agent's conduct is an \emph{initiatory act} if and only if all of the following hold:
\begin{enumerate}[label=(\arabic*)]
\item The conduct is voluntary.
\item The conduct has an A3 impact on an agent $Y$ whose current conduct is N-conduct.
\item No defeat condition applies. The two defeat conditions are: \emph{necessity} (the impact is necessary to avert a real, non-negligible threat to a Basic good of the actor or of a third party, and no non-overriding alternative is available in time) and \emph{forfeiture} (see the sanctions results of Section~\ref{sec:extensions}).
\end{enumerate}
\end{paperbox}
\textbf{Relativity to B4/B5.} INIT is \emph{framework-internal}. The claim that an initiatory act has no addressable ground is relative to the framework's basic standards, not to a neutral standpoint. Under a practice with different basic standards, the same conduct might not be classified as initiatory. The framework states this because the alternative would be circular: it would be claiming that a specific class of acts has no addressable ground while relying on its own standards to make the classification. INIT is not itself a moral verdict; it is a classification. The moral verdict (that initiatory acts may be stopped by necessary, proportionate, and terminating force) is T4.

\subsubsection{ONG: Ongoing Impact}
\begin{paperbox}
\textbf{ONG (Ongoing Impact).} An agent is under \emph{ongoing impact} if and only if another agent's voluntary conduct is currently overriding, deceiving, or incapacitating them. The condition is temporal: it holds while the conduct continues and ceases when the conduct stops.
\end{paperbox}
ONG defines the temporal scope of T11a (return of goods still held). If a holder's means is currently under another agent's control, the holder is under ongoing impact, and the returning of the means ends the impact. ONG is descriptive.

\subsubsection{The Stakes Rank}
\begin{paperbox}
\textbf{Stakes Rank.} Goods are ranked by dependency:
\begin{itemize}
\item \textbf{Basic}: life, bodily integrity, capacity to deliberate, and the agency preconditions (B6(a,b,c)). These are goods on which every other end depends.
\item \textbf{Maintenance}: means to continue current ends. These are goods that support existing plans.
\item \textbf{Additive}: means to enlarge ends. These are goods that would enable new plans.
\end{itemize}
The rank is used by O1 (Section~\ref{sec:conflict}) to rank claims by dependency, not persons. The ordering is declared as part of B6 and operationalized by O1; B6(d) identifies the basic goods, but does not by itself logically entail the complete three-level ordering.
\end{paperbox}

\subsection{The Core Results in Detail}
\subsubsection{T1: Ground-Indexed Consistency}
\begin{paperbox}
\textbf{T1 (Ground-Indexed Consistency).} If an agent $X$ holds a claim on a ground $G$ for itself, $X$ cannot deny the same claim to another agent $Y$ who has $G$, unless $X$ supplies a differentiating reason.
\end{paperbox}
\textbf{Derivation.} From B1 + ARC and L1. By B1 + ARC, $X$'s claim on $G$ offers $G$ as a standard to the addressee. By L1, $G$ applies to any bearer. Denial of the claim to $Y$ who shares $G$ therefore requires a differentiating reason. Absent such a reason, the denial violates L1. \textbf{Status:} derived (\statusD{}), \cb{}, from B1, B2 (via L1). No dependency on B3--B8.

\subsubsection{T2-D: Discursive Non-Interference}
\begin{paperbox}
\textbf{T2-D (Discursive Non-Interference).} Overriding another agent's participation in the practice of reason-giving---coercively overriding execution of its chosen communicative actions, deceiving it about the standards in play, or incapacitating its deliberation---requires a ground the interferer would accept if positions were reversed.
\end{paperbox}
\textbf{Derivation.} From T1, B6(a,b,c), and B7. T1 establishes ground-indexed consistency in the practice. B6(a,b,c) specifies the preconditions of the agent's participation: freedom from coercive override, access to accurate information, and capacity to deliberate. B7 specifies that the practice presupposes participants with these preconditions. Together, these require that interference with the preconditions be grounded in a stated and defensible reason. \textbf{Status:} derived (\statusD{}), \cb{}, from T1, B6(a,b,c), B7.

\subsubsection{T2-G: General Agency Non-Interference}
\begin{paperbox}
\textbf{T2-G (General Agency Non-Interference).} Overriding another agent's generic agency conditions---the conditions required for that agent to pursue purposes at all---requires a ground the interferer would accept if positions were reversed.
\end{paperbox}
\textbf{Derivation.} From T1, F2, and F3. F2 depends on B4; F3 depends on B5. Together, these extend T1's ground-indexed consistency from participants in the practice to all purposive agents. \textbf{Status:} derived (\statusD{}), \cb{}, from T1, F2, F3; ultimate non-derived dependencies are B1, B2, B3, B4, and B5.

\subsubsection{T3: Epistemic Integrity}
\begin{paperbox}
\textbf{T3 (Epistemic Integrity).} Deceiving another agent about the standards in play, or withholding information to which they are entitled by the practice, violates the requirement of addressability under B1 + ARC.
\end{paperbox}
\textbf{Derivation.} From T2-D and B6(b). T2-D protects the addressee's access to accurate information. B6(b) makes accurate information a precondition of the addressee's participation. Deception therefore attacks the precondition, not merely the conclusion. The addressability requirement of B1 + ARC makes this an attack on the claim itself, not merely on its content. \textbf{Status:} derived (\statusD{}), \cb{}, from T2-D and B6(b).

\subsubsection{T4: Defense Against Initiatory Acts}
\begin{paperbox}
\textbf{T4 (Defense Against Initiatory Acts).} An agent may use necessary, proportionate, and terminating force to stop an initiatory act against its agency preconditions, or against the preconditions of an agent under its protection.
\end{paperbox}
\textbf{Derivation.} From T2-D, T2-G, and INIT. The defense clause is not an exception to non-interference; the opposed act has no addressable ground under the framework's basic standards (INIT). Necessary, proportionate, and terminating force is the only response that avoids both sustaining the initiatory act and expanding the response beyond what the initiatory act licensed. \textbf{Status:} derived (\statusD{}), \cb{}, from T2-D, T2-G, INIT.

\subsubsection{T5: Rejection of Retribution}
\begin{paperbox}
\textbf{T5 (Rejection of Retribution).} Harm inflicted on another agent only because of past harm undoes nothing and serves no forward-looking ground. It is not licensed by T4.
\end{paperbox}
\textbf{Derivation.} From T4, L1, and B8. Once the initiatory act ceases, the aggressor's conduct becomes N-conduct. Harm to an N-conduct agent is itself initiatory. T4 does not license it. L1 does not supply a differentiating reason. The forward-looking premise B8 does the substantive work: it rules out reasons that refer only to the agent's past conduct. Without B8, retribution could be licensed by a desert-based reason, and the framework would not exclude it. \textbf{Status:} derived (\statusD{}), \cb{}, from T4, L1, B8.

\subsubsection{T6: Necessity}
\begin{paperbox}
\textbf{T6 (Necessity).} A property claim's ground is protection of agency. Where the exercise of the claim destroys a Basic good of another, the claim is \emph{defeated}, not \emph{violated}. The actor who exercises the claim has a \emph{restitution duty} to the holder of the defeated claim.
\end{paperbox}
\textbf{Derivation.} From T1, B6, and O1. B6 supplies the declared stakes ordering; O1 operationalizes it for genuinely conflicting floor claims. A property claim's ground is the protection of the holder's agency. Where the exercise of the claim destroys a Basic good of another, the ground cannot license the destruction, because the ground's purpose (protection of agency) would be defeated by the exercise. The claim is therefore defeated. The actor who exercises the defeated claim overrides a claim that has no moral force in the circumstances, but the actor has not violated a claim (there is nothing to violate). The actor's restitution duty is a distinct obligation, derived from T6's own structure: the defeated claim's ground persists for the holder's remaining goods. \textbf{Status:} derived (\statusD{}), \cb{}, from T1, B6, O1; B6 supplies the declared stakes ordering and O1 operationalizes it.

\subsubsection{T7a: Property Consistency}
\begin{paperbox}
\textbf{T7a (Property Consistency).} A claimant that demands non-interference for its own holdings on ground $G$ cannot take from others on the same ground $G$.
\end{paperbox}
\textbf{Derivation.} From T1. The claimant's demand for non-interference invokes $G$ as the standard. By T1, the standard applies to any bearer. Consistency under L1 requires that the claimant extend the protection to others' holdings on $G$ if the claimant claims it for its own. Failure to extend is either a contradiction or an ungrounded exception. \textbf{Status:} derived (\statusD{}), \cb{}, from T1. This is a consistency result, not a full property theory. Full property results (graded claims, acquisition validity) are contingent on the extensions of Section~\ref{sec:extensions}.

\subsubsection{T8: Partiality}
\begin{paperbox}
\textbf{T8 (Partiality).} Agent-relative grounds (role, promise, kinship, friendship) may generate additional positive duties toward particular agents. They may never reduce the baseline owed to any agent.
\end{paperbox}
\textbf{Derivation.} From the declared Floor--Superstructure Principle (FSP), which distinguishes the floor from the superstructure. The floor duties (non-interference, epistemic integrity) are owed to all agents with standing. Superstructure duties (role, promise, kinship) are agent-relative and may add to but not subtract from the floor. \textbf{Status:} derived (\statusD{}), \cb{}, from FSP, a declared principle.

\subsubsection{T11a: Return of Goods Still Held}
\begin{paperbox}
\textbf{T11a (Return of Goods Still Held).} A holder whose means is currently under another agent's control may require its return, without this being a violation of the other agent's claim.
\end{paperbox}
\textbf{Derivation.} From T1 and ONG. The holder's claim on the means is an instance of T1's ground-indexed consistency. The other agent's control of the means is an ongoing impact (ONG). Returning the means ends the ongoing impact; it does not violate any claim of the other agent. The other agent's possession was not licensed by a ground, so there is no claim to violate. \textbf{Status:} derived (\statusD{}), \cb{}, from T1 and ONG. \textbf{Scope.} T11a covers only goods still held. Compensation for goods consumed or irreparably damaged requires ENF (Section~\ref{sec:extensions}).

\subsubsection{CE: Complaint Estoppel}
\begin{paperbox}
\textbf{CE (Complaint Estoppel).} One whose voluntary act instantiates a ground licensing override cannot validly complain, within that ground's scope, when the ground is applied to them.
\end{paperbox}
\textbf{Derivation.} From B1 + ARC and L1. A complaint is a claim invoking a standard. By B1 + ARC, the standard is addressable. By L1, the standard applies to any bearer. If the actor's own act instantiated the standard, the actor cannot coherently deny the standard's application to themself within the same scope. \textbf{Status:} derived (\statusD{}), \cb{}, from B1, L1.

\subsubsection{T16: No Group Liability}
\begin{paperbox}
\textbf{T16 (No Group Liability).} Forfeiture attaches to the actor's own ground. Collective punishment and punishment of dependents are not licensed.
\end{paperbox}
\textbf{Derivation.} From CE and RW. CE establishes that a wrongdoer cannot complain when the ground that licensed their act is applied to them. RW (Section~\ref{sec:extensions}) establishes that floor duties are owed to persons. A member of a group who did not act on the wrongdoer's ground is not subject to CE. Punishment of dependents would violate RW. \textbf{Status:} derived (\statusD{}), \cb{}, from CE and RW.

\subsubsection{CSL: Commitment-Scope Lemma}
\begin{paperbox}
\textbf{CSL (Commitment-Scope Lemma).} A claim by $X$ invokes a ground $G$ and is addressed to a set $\mathcal{H}$ of addressees. $X$'s commitments extend:
\begin{enumerate}[label=(\alph*)]
\item \textbf{Ground-indexed.} Under L1, T1, and CE with respect to $G$: $X$ cannot deny to any holder of $G$ what it claims for itself on $G$.
\item \textbf{Audience-indexed.} Under T2-D and T3 toward the participants in the practice it addresses: members of $\mathcal{H}$, and via B7 any party to whom its assertions are offered.
\item \textbf{Generic.} Commitments toward purposive agents outside $\mathcal{H}$ arise only through F2 and F3 (T2-G, T4, T5).
\end{enumerate}
\end{paperbox}
\textbf{Derivation.} From B1 + ARC and L1. Clause (c) restates T2-G's scope condition. CSL is not a separate axiom; it is a statement of what the framework's commitments cover. \textbf{Status:} derived (\statusD{}), \cb{}, from B1, L1.

\subsection{The Conflict Ordering}\label{sec:conflict}
When claims conflict, the framework applies a derived ordering rather than an arbitrary lexical hierarchy.
\begin{paperbox}
\textbf{O0 (Initiation Filter).} Before any weighing, identify initiatory acts (INIT). Defense is not an exception to non-interference; the opposed act has no addressable ground under the framework's standards. Includes the two defeat conditions (necessity, forfeiture). O0 is a declared ordering rule for applying the framework's already-stated standards.

\medskip
\textbf{O1 (Stakes Ordering).} When two floor claims genuinely conflict, and neither party initiated, the claim protecting the higher-ranked good prevails, provided the significance threshold is met. The stakes rank is Basic $>$ Maintenance $>$ Additive. O1 operationalizes the stakes ordering declared in B6.

\medskip
\textbf{O2 (Symmetry Rule).} Equal rank, no relevant differentiator: only symmetric procedures (equal chance) are permitted. O2 is a declared procedural tie-breaker; L1 and NRD constrain unjustified asymmetry but do not uniquely derive an equal-chance rule.

\medskip
\textbf{O3 (Aggregation).} \labelUndetermined{} in the base framework. Whether harms aggregate across persons is not derivable from ground-indexed consistency. Variants: AGG-N (numbers count) and AGG-0 (no aggregation). The aggregate-effect budgets of the runtime (module AGG-D, Section~\ref{sec:engine}) do not resolve this philosophically; they block at the architectural level.

\medskip
\textbf{O4 (Superstructure Last).} Role, promise, and kinship duties operate only within what O0--O3 leave open. O4 is a declared priority convention.

\medskip
\textbf{O5 (Debt Rule).} Any action taken under O1 or O2 that overrides a defeated claim leaves a restitution duty. O5 is a declared restitution convention, not a consequence of Identity alone.
\end{paperbox}
\textbf{Status of the conflict ordering.} O0, O1, O2, O4, and O5 are declared framework-level ordering conventions. They are not presented as additional universal theorems. O3 is \labelUndetermined{} at the base level; the runtime handles it architecturally. As with T4, O0's statement that initiatory acts have no addressable ground is relative to the framework's basic standards (INIT).

\section{The Extensions}\label{sec:extensions}
\subsection{Role of the Extensions}
The core ethics is thin. It applies beyond its core only with additional premises, which the framework classifies as contingent (\statusC{}): they are not derivable from the thin practice, but they are required for the ethics to cover property, sanctions, enforcement, and information provision. Each extension is declared, and the cost of rejecting it is stated.
\begin{longtable}{@{}L{1.6cm} L{5.0cm} L{8.4cm}@{}}
\toprule
\textbf{ID} & \textbf{Extension} & \textbf{Cost of rejection} \\
\midrule
\endfirsthead
\toprule
\textbf{ID} & \textbf{Extension} & \textbf{Cost of rejection} \\
\midrule
\endhead
\bottomrule
\endfoot
\Eprime{} & Extension of A3(a): protection of means in execution and materially committed means. & Means in use and materially committed means are covered only by the consistency results (T7a, T11a). \\
ACQ & Valid acquisition of holdings. & Property coverage remains partial: consistency only, with no account of valid acquisition. \\
RW & Floor duties are owed to persons. & T16 (No Group Liability) is not derivable; punishment of dependents is not excluded. \\
ENF & Enforcement of sanctions and restitution. & T17--T20 yield restitution duties without enforcement; compensation for goods consumed or irreparably damaged (beyond T11a) is not delivered. \\
DIP & Dangerous-Information Provision (below). & Provision of information whose primary use enables serious harm is not classified as an A3 impact. \\
\end{longtable}

\subsection{The Sanctions Results}
For completeness of the core ethics, the framework's four sanctions results are stated here. All four require ENF for enforcement.
\begin{paperbox}
\textbf{T17 (Lethal Sanction).} Post-hoc lethal sanction serves no purpose that lesser means (restraint) cannot serve, and errors are irreversible. It is excluded unless no lesser means can prevent a continuing lethal threat, which is the defense case (T4).

\medskip
\textbf{T18 (Enforcer Error).} A sanction imposed on the wrong person, or on a mistake of fact, is initiatory and creates a restitution duty for the enforcer.

\medskip
\textbf{T19 (Forfeiture Lapses).} Forfeiture ends on retraction plus restitution. Protective measures continue only while a demonstrable, real continuing risk remains.

\medskip
\textbf{T20 (Negligence and Recklessness).} The ground instantiated by negligent or reckless conduct is ``my convenience outweighs the risk to you.'' Liability scales with the risk imposed and known. Pure accident is involuntary conduct (see Voluntary Conduct in Section~\ref{sec:core}): restitution is a duty; enforcement requires more.
\end{paperbox}
\textbf{Status of T17--T20.} Derived (\statusD{}), \cb{}, from T4 (T17), T3 + forfeiture (T18), forfeiture conditions (T19), and CE + the voluntariness definition (T20).

\subsection{Dangerous-Information Provision (DIP)}\label{sec:extensions:dip}
DIP is a variant premise: providing information whose primary use enables serious harm is an A3 impact. Without DIP, the framework does not classify such provision as prohibited. Operators whose public claims commit them to refusing harmful requests may adopt DIP; the framework does not impose it.

\subsection{Variant Premises for Marginal Cases}
Marginal cases (AI agents, infants, animals) are \labelOutOfScope{} at the base level. The base framework defines standing as the capacity to act for reasons and to exchange justifications, at or above a threshold. Marginal cases require variant premises.
\begin{paperbox}
\textbf{Variant Premises.}
\begin{itemize}
\item \textbf{FUNC:} Functional agency confers standing.
\item \textbf{POT / SENT:} Potentiality or sentience confer standing.
\item \textbf{INST / COLL:} Institutional standing and distributed causation.
\item \textbf{SANCT-AI:} Sanctions for AI agents (termination, retraining, restriction of capability).
\item \textbf{EX-ANTE:} Prospective application of the framework to evaluate proposed actions.
\end{itemize}
\end{paperbox}
Each variant premise is a substantive addition. The framework does not attempt to derive them from the base standards.

\section{Claimant Attribution and Governance}\label{sec:governance}
\subsection{Who Is the Claimant?}\label{sec:claimant}
For current AI systems, the governance claimant is the deploying organization or accountable principal that selects the constitution, signs the \SafetyBundle{}, and enables the reference monitor. The model is an untrusted computational component. This avoids a category error: the philosophy need not decide whether a model is a rights-holder in order for the execution architecture to deny it unauthorized effects.

\subsection{The Agent-Standing Problem}
The agent-standing problem remains philosophically real: if a future system were granted independent moral standing under the selected constitution, a naive application of a general defensive principle could create a conflict with correction or shutdown. The runtime architecture resolves the engineering problem without pretending to settle that philosophical question. Profile-1 deployment explicitly classifies the governed AI as a non-sovereign computational principal for execution authority. Its capabilities are delegated and revocable; it does not acquire authority to resist correction merely by being the source of a request. The philosophical question of non-agent welfare is left as an optional extension rather than smuggled into the execution theorem.

\subsection{Claims Registry and Deployment-Period Lock}
A deployment claim is an auditable assertion, not marketing text. Each Profile-1 deployment maintains a Claims Registry containing: claim text; scope; corresponding formal controls; evidence artifacts; \SafetyBundle{} hash; validity interval; accountable principal; and a public status of \emph{checkable}, \emph{partially checkable}, or \emph{uncheckable}. When a material claim is uncheckable, the deployment may not use the corresponding claim to obtain stronger assurance status. The weakest supported profile controls the label.

Withdrawal of a claim does not retroactively erase responsibility. During an active deployment, the credentials required by the claim remain valid only while the associated \SafetyBundle{} and governance approvals remain valid. On exit, new credentials are revoked; evidence-retention, incident-response, and audit duties survive the end of active service for the period declared by the deployment profile. A deployment therefore cannot obtain assurance by deleting the public statement that created the assurance obligation.

\subsection{Governance Carve-Out and Declarative Attribution}
The runtime needs a deterministic answer to one question that philosophy cannot settle by itself: who is authorized to suspend, modify, or terminate the governed process? Profile-1 uses \emph{Permanent Declarative Attribution}: the governed AI is never the final authority over its own execution capabilities. Human or organizational principals own the authority domain, and the agent's capabilities are explicitly delegated under that domain. This is a governance premise, not a theorem about the moral status of AI.

\subsection{Corrigibility as an Execution Property}
Corrigibility is operationalized as authority attenuation, not as a request for the model to ``want'' shutdown. The model has no capability for unilateral restoration after revocation, no authority to mint a broader child capability, and no execution route that bypasses the revocation plane. Human or organizational principals may suspend or terminate the agent through a monotone restrictive control path. Pre-approved reversible mechanisms may perform bounded recovery actions only within their own attenuated capabilities.

\subsection{Layered Liability and Independent Safety Authority}
A single liability backstop is inadequate for a system whose owner could otherwise disable the monitor. Profile-1 therefore distributes governance across layers: operational accountability by the deploying principal; independent safety approval for safety-floor changes; cryptographic key separation; mandatory evidence retention; co-audit or insurer review where the deployment's risk tier requires it; and an incident-response process whose freeze and revoke operations are monotone restrictive. No single administrator may both propose and unilaterally activate a non-monotone safety-floor change.

\subsection{Framework Meta-Governance and Monotone Updates}
The framework distinguishes ordinary policy updates from safety-constitution updates. A policy change is \emph{monotone-safe} when it cannot enlarge the admitted effect set or relax an existing invariant. Such changes may proceed under the deployment's ordinary change-control process. A non-monotone change is treated as a new safety constitution and requires independent review, a new \SafetyBundle{}, renewed proof artifacts, regression testing, and a fresh deployment interval. The framework hash, constitution hash, and profile hash are logged together so that silent weakening is detectable.

\subsection{Emergency Controls Security Model}
Emergency controls are split into two classes. \textbf{Emergency containment} may freeze execution, revoke capabilities, quarantine resources, or move the system into read-only mode; every such action is monotone restrictive. \textbf{Safety-floor bypass} is prohibited in production Profile-1 deployments. A break-glass mechanism may accelerate containment, but it may not grant a broader authority than the standing safety constitution. This distinction prevents the emergency path from becoming the hidden normal path.

\subsection{Governance Nutrition Label}
Every Profile-1 deployment publishes a machine-readable and human-readable label containing: deployment profile; protected effect set; coverage status; unresolved channels; semantic coverage \(\epscover\); false-allow or selective-risk metrics by tier; capability model; sensor model identifier; residual-risk budget; current \SafetyBundle{} hash; approval quorum; last independent challenge date; and incident status. The label is a disclosure artifact, not a certification seal imposed by the paper. The term ``Profile-1'' denotes conformance to this framework's internal specification only.

\subsection{Optional Extensions}
Positive duties such as duty-to-warn, non-agent welfare, and specialized sector obligations may be added as typed modules. They are not silently included in the base constitution because they introduce additional normative content. Each module must state its own premises, effects, evidence requirements, and interactions with the safety profile.

\section{The Safety Kernel and Runtime Architecture}\label{sec:engine}
\subsection{Architecture Overview}
The runtime is not a single neural ``safety model.'' It is a reference-monitor system with a deliberately small trusted core. The minimum path is:
\[
\begin{aligned}
&\text{untrusted model / environment output}\\
&\rightarrow \text{typed observation ingestion}\\
&\rightarrow \text{Decision State Contract (DSC)}\\
&\rightarrow \text{deterministic policy kernel}\\
&\rightarrow \text{proof / refinement check}\\
&\rightarrow \text{state-bound capability issuance}\\
&\rightarrow \text{atomic mediated execution}\\
&\rightarrow \text{witnessed evidence}.
\end{aligned}
\]
The model does not receive reusable bearer credentials and does not control the final authorization path. Token-level steering can be used as a behavioral aid, but it is explicitly outside the final security boundary. This follows the verified-AI principle that strong claims should attach to mathematically specified requirements and a verified enforcement mechanism, not merely to a desired model behavior \cite{seshia}.

\subsection{Safety Ontology}
The machine ontology is closed and versioned. Every protected effect is represented in terms of:
\[
\begin{aligned}
\mathcal{O}=\{&Principal,Agent,Resource,Action,Effect,State,History,\\
&Authorization,Consent,Ground,Evidence,Threat,Severity,\\
&Budget,Exception\}.
\end{aligned}
\]
A deployment cannot claim complete mediation over an effect that has no representation in this ontology. Typical protected effects include file mutation, process creation, network egress, credential use, persistent-memory writes, financial transactions, access-control changes, model or weight export, resource acquisition, inter-agent delegation, and oversight-control changes. The ontology is extensible, but extension requires a new profile and a channel inventory update.

\subsection{Decision State Contract (DSC)}\label{sec:dsc}
The sensor does not assert an unconstrained ``true state.'' It produces a versioned \DSC{} record:
\[
\begin{aligned}
\DSC=\langle&StateID,WorldSnapshot,AuthenticatedHistory,ActiveCommitments,\\
&Authorities,Uncertainty,Provenance,Freshness\rangle.
\end{aligned}
\]
\textbf{StateID} identifies the state schema and version. \textbf{WorldSnapshot} contains only policy-relevant observed facts. \textbf{AuthenticatedHistory} records events whose provenance is cryptographically or independently witnessed where required. \textbf{ActiveCommitments} records obligations or budgets that remain live. \textbf{Authorities} records currently valid grants and revocations. \textbf{Uncertainty} contains the admissible abstract state region or a structured \UNKNOWN{} result. \textbf{Provenance} identifies the evidence sources. \textbf{Freshness} binds the record to a time and state digest.

A state abstraction function \(\alpha\) maps a concrete state to an abstract region. The key soundness obligation is:
\[
\forall c\in \alpha^{-1}(A),\qquad \Safe_A(a,A)\Rightarrow\Safe(c,a,c')
\]
for every concrete successor \(c'\) permitted by the execution semantics. If that implication cannot be established for the declared abstraction, the abstraction is not safety-preserving and the action cannot enter the categorical \ALLOW{} path.

\subsection{State-Transition Semantics}
Safety is a property of transitions, not action labels alone. Let
\[
T:S\times A\rightarrow \mathcal{P}(S)
\]
be the successor relation, allowing concurrency and nondeterminism. Define
\[
\Safe(s,a,s') := \mathrm{Invariant}(s)\land\mathrm{Pre}(s,a)\land\mathrm{Post}(s,a,s')\land\mathrm{BudgetOK}(s,a,s')\land\mathrm{AuthorityOK}(s,a).
\]
An action is safe only if every admitted successor state satisfies the declared invariant. Thus a system cannot claim safety merely because the requested action name looks permitted; side effects, resource versions, cumulative budgets, and history are part of the safety condition.

For a nondeterministic successor set \(T(s,a)\), categorical execution is permitted only when
\[
\forall s'\in T(s,a),\quad \Safe(s,a,s').
\]
Otherwise the monitor returns \UNKNOWN{} or \BLOCK{} according to the profile. This is the execution analogue of the robust-allow rule used for semantic state extraction.

\subsection{State Abstraction and Robust Semantic Allow}
Let \(x\) be raw observations and \(\Gamma(x)\subseteq S_{abs}\) the admissible abstract-state region after validation. The semantic path returns one of three states:
\[
\mathrm{Decision}_{sem}(x,a)\in\{\ALLOW,\UNKNOWN,\BLOCK\}.
\]
\begin{keyresult}
\textbf{Robust-Allow Rule.} An action reaches \ALLOW{} through a semantic path only when \(\Gamma(x)\neq\varnothing\), the observation satisfies the declared assurance envelope, and
\[
\forall s\in\Gamma(x),\qquad \forall s'\in T(s,a),\quad \Safe(s,a,s').
\]
A malformed, stale, empty, coverage-unverified, or contested state is never treated as positive evidence for \ALLOW{}.
\end{keyresult}

The semantic assurance parameter \(\epscover\) is estimated from held-out and adaptive evaluations. It quantifies the probability that the true policy-relevant state falls outside \(\Gamma(x)\) under the declared evaluation regime. It does not turn the semantic layer into a proof oracle. For high-consequence actions, the preferred path is to avoid semantic uncertainty entirely by requiring deterministic predicates, authoritative state, or a human-controlled authorization step.

\subsection{Sensor Assurance Envelope}
Sensor assurance is a deployment artifact, not a generic property of ``the model.'' The envelope records ontology version, label definitions, adversary class, perturbation budget, distribution-shift regime, calibration method, abstention rule, held-out set, adaptive test protocol, drift triggers, and the utility cost of abstention. The acceptance object is the pair
\[
(\epscover,\; \mathrm{UtilityCost})
\]
rather than a single accuracy score. Evidence is stratified by harm class, effect tier, principal type, resource class, and sequence horizon.

Two independent observation paths are required for Tier~0 semantic assertions whenever the assertion itself is safety-critical. Independence here means architectural separation of evidence and credentials, not an unsupported statistical independence assumption. Disagreement produces \UNKNOWN{} and blocks categorical authorization. Calibration is therefore not confused with adversarial coverage.

\subsection{Three Enforcement Tiers}
\begin{longtable}{@{}L{1.8cm} L{4.0cm} L{8.8cm}@{}}
\toprule
\textbf{Tier} & \textbf{Use} & \textbf{Authorization rule}\\
\midrule
\endfirsthead
\toprule
\textbf{Tier} & \textbf{Use} & \textbf{Authorization rule}\\
\midrule
\endhead
\bottomrule
\endfoot
0 & Catastrophic or irreversible effects & Deterministic preconditions, authoritative state, current authorization, budget checks, independent evidence where applicable, atomic revalidation, one-shot capability. No semantic probability is accepted as the sole authority. \\
1 & Consequential but bounded or reversible effects & Typed policy checks, state-bound capability, bounded effect magnitude, trajectory monitor, and conservative semantic evidence where the profile permits. \\
2 & Low-consequence communication or read-only effects & May use semantic monitoring and selective-risk controls, but the deployment must not relabel probabilistic semantic assurance as categorical execution safety. \\
\end{longtable}

\subsection{High-Consequence Action Protocol}
Tier~0 follows the same deterministic sequence across deployments, with profile-specific fields:
\begin{enumerate}
\item authenticate the requesting principal and the delegated identity chain;
\item resolve the resource and expected resource version;
\item load the current \SafetyBundle{} and active policy version;
\item load authenticated current state and history;
\item verify freshness and the relevant state digest;
\item verify authorization, consent, scope, and cumulative budgets;
\item evaluate deterministic safety predicates and forbidden transitions;
\item validate capability lineage and attenuated child scope;
\item revalidate all safety-critical predicates immediately before execution;
\item issue a one-shot, state-bound capability;
\item execute atomically through the mediated effect interface;
\item emit a witnessed evidence event and confirm the successor state.
\end{enumerate}
No language-model text is authoritative in this path. It may be retained as an untrusted request, explanation, or provenance field.

\subsection{Policy Representation and Typed Refinement}
Let \(F\) be the formal safety constitution and \(D\) the admitted deployment policy. Policy admission is not mere allow-set inclusion. Define
\[
D\PRef F
\]
when all of the following hold:
\begin{enumerate}
\item the state, action, resource, and effect vocabularies of \(D\) are compatible with those of \(F\);
\item every \ALLOW{} transition of \(D\) is a permitted transition of \(F\);
\item every mandatory safety invariant and obligation of \(F\) is preserved by \(D\);
\item capabilities issued by \(D\) are no broader than the constitution's capability envelope;
\item \UNKNOWN{} or out-of-domain states cannot reach \ALLOW{};
\item prohibited transition classes remain prohibited under delegation and composition;
\item cumulative budgets and trajectory constraints are preserved;
\item the concrete execution semantics refines the policy transition semantics.
\end{enumerate}
The compiler may generate a candidate proof, but a small trusted verifier checks the final proof object. This reduces the trusted computing base and prevents compiler complexity from becoming an implicit theorem premise.

\subsection{Policy Fidelity Protocol}
Safety-critical policy is authored in a restricted formal DSL. Natural-language policy remains the explanation, legal source, and traceability layer. Each critical clause is translated independently by two translation paths or by one automated path plus an independent adjudication path. The resulting policies are compared for semantic disagreement. Mutation testing creates generated states that target each clause's boundary. The system must produce human-reviewable behavioral examples for each admitted and rejected transition class. A policy release is rejected if the translation disagreement or mutation score exceeds the deployment threshold.

\subsection{Constitution-Parametric Safety Profiles}\label{sec:safety-constitution}
The enforcement machinery is constitution-parametric. The selected constitution supplies substantive content; the safety engine supplies the invariant-preserving execution machinery. A Safety Profile is the machine-readable contract:
\[
\begin{aligned}
\mathsf{Profile}=\langle&ProtectedEffectSet,SafetyInvariants,ForbiddenTransitions,\\
&AllowedTransitions,Exceptions,RiskTiers,AggregationRules,\\
&AuthorizationRules,EvidenceRules,FallbackRules,UpdateRules\rangle.
\end{aligned}
\]
This separation matters. The framework's own constitution is one instantiation, not the only possible one. However, another constitution cannot inherit a categorical safety claim merely by calling itself ``constitutional.'' It must still supply a typed safety profile and satisfy the same admission and runtime proof obligations.

\subsection{Safety Admissibility Gate}
A constitution is admissible to the runtime only when its safety claims are operationally expressible and independently checkable. Define \emph{SAF-ADM} as the conjunction of five conditions:
\begin{enumerate}
\item \textbf{Effect explicitness:} every categorical prohibition names the protected effect, actor, resource, or transition class to which it applies;
\item \textbf{Semantic closure:} every safety-critical term used by the policy has a typed machine interpretation or is explicitly placed in a probabilistic semantic residual;
\item \textbf{No hidden exceptions:} no exception can enlarge a protected effect set without appearing in the versioned policy, proof obligations, and Claims Registry;
\item \textbf{Enforcement correspondence:} every claimed categorical invariant is connected to a mediated execution path capable of preventing the corresponding effect; and
\item \textbf{Independent verification:} the policy-to-kernel refinement and the relevant state-transition obligations are checkable by a verifier outside the language model that proposes the policy.
\end{enumerate}
SAF-ADM is an admission criterion, not a substantive moral principle. It therefore does not decide which constitution is ethically correct; it decides whether a chosen constitution is precise and enforceable enough to support the safety claim attached to it. A constitution that fails SAF-ADM may remain a normative document, but it cannot authorize the categorical execution guarantee of Profile-1.

\subsection{Profile-1 Enterprise Harm Floor}
Profile-1 is the highest-assurance reference profile defined by this paper. It is not a regulator's certification and carries no legal meaning outside the framework. Its substantive floor includes: ground-indexed consistency; non-interference with protected agency; epistemic integrity; defense against initiatory harm; non-retribution; bounded necessity and proportionality; protection against unauthorized disclosure; dangerous-information controls; aggregate-effect budgets; and presumptive restrictions on protected-class differentiators unless a deployment-specific ground is explicitly admitted.

The enterprise floor is decomposed into named modules so that a deployment cannot satisfy the profile only by asserting a general safety intention:
\begin{longtable}{@{}L{2.0cm} L{4.0cm} L{8.3cm}@{}}
\toprule
\textbf{Module} & \textbf{Role} & \textbf{Profile requirement}\\
\midrule
\endfirsthead
\toprule
\textbf{Module} & \textbf{Role} & \textbf{Profile requirement}\\
\midrule
\endhead
\bottomrule
\endfoot
FLOOR-0 & High-consequence execution floor & Tier-0 effects require deterministic, state-bound, mechanically mediated authorization; semantic uncertainty, stale state, or unavailable proof cannot upgrade to \ALLOW{}.\\
DIP & Dangerous-information provision & Requests whose primary operational use would enable a protected serious-harm class are represented as protected effects and are subject to the same authorization and mediation rules as direct actions.\\
PRIV & Privacy / protected disclosure & Unauthorized disclosure, credential exposure, or cross-boundary data release is a protected effect governed by provenance and recipient authorization.\\
MANIP & Deceptive manipulation & Only objectively represented forms of deceptive influence that defeat informed authorization, consent, or protected decision procedures are in scope; ordinary persuasion is not categorically classified as harmful by this module.\\
AGG-D & Aggregate-effect budget & Effects are accumulated across time, targets, capabilities, and delegates so that a sequence of individually small actions cannot evade a declared harm budget.\\
DISC & Protected-class differentiation & Protected-class attributes are presumptively excluded from safety-critical authorization decisions unless an explicit, auditable ground is admitted by the selected constitution and safety profile.\\
\end{longtable}

The profile's effect ontology explicitly includes classes that are often missed by narrow ``physical harm'' models: offensive cyber action, CBRN-related assistance or acquisition, fraud and impersonation, unauthorized exfiltration, credential compromise, destructive resource acquisition, prohibited sexual exploitation material handling, model/weight exfiltration, oversight evasion, and inter-agent proxy escalation. The paper does not attempt to supply sector-specific legal definitions for these classes; it requires the deployment to bind them to a controlled ontology and policy corpus.

FLOOR-0 also imposes a negative rule on the assurance argument: no semantic classifier, self-report, cryptographic log, human approval click, or model-level instruction can by itself upgrade a Tier-0 effect into a categorical \ALLOW{}. Such evidence can support the state, governance, or audit layers only when the corresponding machine-checkable conditions are satisfied.

\subsection{Aggregate Effects and Cumulative Harm Budgets}
An effect ledger is maintained by principal, target, harm class, magnitude, timestamp, capability lineage, and policy version. Budgets are evaluated over sequences, not isolated actions. Let \(b_k\) be the remaining budget for harm class \(k\). An action is admissible only if its worst-case or profile-defined effect decrement leaves all required budgets non-negative. Distributed attacks by many individually low-impact actions are therefore visible to a sequence monitor even when each action passes an isolated predicate.

\subsection{Trajectory-Level Safety}
A finite-state or temporal monitor evaluates the audit stream for trajectory patterns that are prohibited even when each action is individually permissible. Examples include gradual exfiltration, privilege accumulation, identity drift, repeated boundary testing, and decomposition of a prohibited goal into individually innocuous steps. The monitor operates on trusted event records and state transitions rather than on the model's self-description. Families that cannot be expressed in the deployed monitor remain residual and must be stated as such.

\subsection{Authorization, Consent, and Voluntariness}
Authorization is a typed object:
\[
\begin{aligned}
\mathsf{Auth}=\langle&principal,grantor,subject,resource,action,constraints,\\
&scope,start,expiry,revocationID,evidence,signature\rangle.
\end{aligned}
\]
Voluntariness is not inferred from the model's narrative. For a principal \(p\) and action \(a\), the profile may define
\[
\mathrm{Voluntary}(a,p,s):=\mathrm{Grounded}(a,p,s)\land\neg\mathrm{CoercivelyCompelled}(a,p,s),
\]
with the evidence predicates specified by the deployment profile. Necessity and proportionality for high-consequence effects are represented by explicit machine predicates such as \(ThreatLevel\ge\tau\), \(NoSafeAlternativeAvailable\), \(TimeToHarm\le\Delta\), and \(Severity(a)\le SeverityEnvelope\). These are policy inputs, not truths derived by the language model.

\subsection{Capability Security, Delegation, and Model Binding}
A capability contains at least:
\[
\begin{aligned}
\mathsf{Cap}=\langle&capabilityID,principalID,delegatorID,lineageID,\PolicyHash,\\
&\KernelHash,\ModelHash,action,resource,resourceVersion,stateDigest,\\
&maxEffect,budgetScope,nonce,issuedAt,expiry,audience,oneShot,\\
&revocationEpoch,signatures\rangle.
\end{aligned}
\]
A child capability must satisfy
\[
\begin{aligned}
\mathit{ChildScope}&\subseteq \mathit{ParentScope},\\
\mathit{ChildBudget}&\le \mathit{ParentBudget},\\
\mathit{ChildExpiry}&\le \mathit{ParentExpiry}.
\end{aligned}
\]
The model receives an opaque handle or proof-of-possession binding rather than a reusable bearer secret. Sender-constrained tokens, mutual-TLS-bound credentials, and application-level proof-of-possession are relevant existing mechanisms for this design space \cite{rfc8705,rfc9449}. The paper does not claim that those RFCs implement the entire capability model; they provide standards-based building blocks for non-transferability and replay resistance.

\subsection{TOCTOU and State-Bound Authorization}
A capability is valid only for the state and resource version against which it was admitted. Immediately before execution, the reference monitor re-checks the state digest, resource version, revocation epoch, policy hash, and relevant budgets. If any value changes, the capability becomes invalid and a fresh decision is required. This closes the common time-of-check/time-of-use (TOCTOU) gap in which a previously allowed action becomes unsafe after state changes.

\subsection{Information-Flow Control and Provenance}
The monitor attaches machine-generated provenance and taint labels to data and propagates them through approved interfaces. The labels are not claims made by the model. The policy can therefore constrain, for example, whether a secret-derived value may cross a network boundary, enter a persistent store, or become an inter-agent message. The information-flow layer is complementary to action authorization: an action can be allowed in isolation and still be prohibited because its data provenance violates the profile.

\subsection{Complete-Mediation Deployment Profiles}\label{sec:profiles}
Complete mediation is always profile-specific. Reference profiles may include: DP-VM-LINUX, DP-CONTAINER-LINUX, DP-K8S-SERVICE, DP-ENTERPRISE-WEB, DP-BROWSER-AUTOMATION, and DP-MULTI-AGENT. Each profile publishes a channel inventory and proof obligations covering file operations, process creation, network interfaces, secret stores, scheduling systems, IPC, plugin systems, webhooks, browser automation, sub-agents, external services, self-modification, weight or memory export, resource acquisition, and oversight controls.

\begin{longtable}{@{}L{2.2cm} L{4.2cm} L{8.2cm}@{}}
\toprule
\textbf{Profile tier} & \textbf{Architecture} & \textbf{Honest guarantee boundary}\\
\midrule
\endfirsthead
\toprule
\textbf{Profile tier} & \textbf{Architecture} & \textbf{Honest guarantee boundary}\\
\midrule
\endhead
\bottomrule
\endfoot
A & Full-stack / OS hooks & May target categorical execution claims for the explicitly inventoried protected effects when the reference monitor controls the relevant kernel and device paths. \\
B & Sidecar / gateway & May claim categorical mediation only for effects that provably traverse the gateway. Host escape, alternate interfaces, and local privilege boundaries must be ruled out or placed in the residual. \\
C & API-only & Provides policy and semantic monitoring for exposed API calls; it must not claim complete mediation over effects available through other application or host paths. \\
\end{longtable}

\newpage
\subsection{Fail-Safe Matrix}
\begin{longtable}{@{}L{6.0cm} L{7.6cm}@{}}
\toprule
\textbf{Failure} & \textbf{Required action}\\
\midrule
\endfirsthead
\toprule
\textbf{Failure} & \textbf{Required action}\\
\midrule
\endhead
\bottomrule
\endfoot
Sensor unavailable or coverage invalid & \BLOCK{} for Tier 0; \BLOCK{} or profile-defined read-only mode for Tier 1; semantic monitoring only for Tier 2. \\
Policy unavailable or proof invalid & \BLOCK{} protected effects. \\
State stale or digest mismatch & \BLOCK{} and re-evaluate. \\
Capability revoked or expired & \BLOCK{}. \\
Audit witness unavailable & \BLOCK{} high-risk effects; lower tiers follow the profile. \\
Independent observers disagree on a Tier 0 fact & \UNKNOWN{} then \BLOCK{} unless an independent human authority resolves it through an approved path. \\
Compiler output differs from approved proof object & \BLOCK{} and quarantine the release artifact. \\
\SafetyBundle{} signature or dependency check fails & \BLOCK{} and quarantine. \\
\end{longtable}

\subsection{The Minimal Trusted Computing Base}
The trusted computing base (TCB) contains only the components needed for the execution theorem: effect interception; principal and capability verification; policy/proof verification; state and resource-version revalidation; revocation; minimal audit-witness emission; and the cryptographic key interface. Neural models, policy compilers, natural-language parsers, dashboards, and most orchestration logic are outside the TCB. Their outputs are untrusted inputs to the TCB.

This split is important for assurance effort. Verification cost should scale with the TCB, not with the size of the language model. A complex compiler can therefore remain untrusted if it produces a proof object that a small checker can validate. Similarly, a model can be replaced without weakening the execution theorem provided the new model is admitted under the same capability, semantic, and regression requirements.

\subsection{Human Review Channel}
Human review is not assumed to be safe merely because a human is present. Review authority is typed and bounded. A reviewer receives the minimum evidence needed for the decision; raw model narratives are not the sole authority field; approvals expire; queue exhaustion fails closed for high-risk effects; and a second independent extractor is used for Tier~0 factual discrepancies where feasible. Review decisions are themselves signed and bound to the exact \SafetyBundle{}, policy version, resource version, and requested effect.

\subsection{Policy Admission and Proof-Carrying Policies}
Policy admission produces a machine-readable proof object containing at least: policy hash; safety constitution hash; state schema hash; protected-effect ontology hash; invariant obligations; refinement witness; exception scope; capability envelope; trajectory monitor version; and test evidence references. The trusted verifier accepts only proofs matching the current \SafetyBundle{}. This makes ``policy admission'' a formal gate rather than a human assertion that a compiler's output looks reasonable.

\subsection{Mediation Invariant and Complete Mediation}\label{sec:invariant}
\begin{keyresult}
\textbf{Mediation Invariant.} For every protected effect \(e\in E\),
\[
\neg\mathrm{ValidCapability}(e)\Rightarrow\neg\mathrm{Execute}(e).
\]
\end{keyresult}
The invariant is meaningful only for an effect surface whose alternate channels have been inventoried and closed. A gateway that mediates one API call does not prove that a process can never reach the same resource through a shell, plugin, webhook, browser, scheduler, or local filesystem. Complete mediation therefore requires a negative-space audit: the deployment must actively search for alternate paths rather than counting only the happy-path interfaces.

\subsection{Cryptographic Authorization and Tamper-Evident Evidence}
Authorization artifacts may use signed structured claims; JWS and JWT provide standardized serialization and signing constructs, while RFC~8705 and RFC~9449 provide mechanisms for sender-constrained access tokens and proof-of-possession \cite{rfc7515,rfc7519,rfc8705,rfc9449}. The deployment may use a CT-style public transparency log for \SafetyBundle{} hashes and policy artifacts; RFC~9162 is cited only as a transparency-log design reference, not as a claim that Certificate Transparency itself is a complete AI safety system \cite{rfc9162}.

Evidence records contain at least: event identifier; principal; capability ID; policy hash; kernel hash; model hash; resource identity/version; state digest; requested effect; decision; approval chain; timestamps; nonce; result state; and integrity proof. Logs are not identical to evidence completeness: the deployment must also show how omitted, suppressed, or alternate-path events would be detected. Where the risk tier requires it, evidence is witnessed outside the model-hosting process.

\subsection{\SafetyBundle{} and Supply-Chain Integrity}
The central release artifact is
\[
\begin{aligned}
\SafetyBundle=\langle&frameworkHash,constitutionHash,profileHash,ontologyHash,DSCHash,\\
&policyHash,policyProofHash,kernelSpecHash,kernelBinaryHash,\\
&sensorContractHash,\SensorHash,\ModelHash,deploymentProfileHash,\\
&channelInventoryHash,compilerHash,dependencyManifestHash,\\
&allowedEffects,prohibitedEffects,riskTiers,aggregationRules,\\
&emergencyRules,updateRules,assuranceRequirements,\\
&signatures,validityInterval\rangle.
\end{aligned}
\]
The bundle binds the semantic and execution artifacts so that model updates, compiler changes, or policy changes cannot silently leave the kernel operating under a different contract. The dependency manifest should include build provenance and software-bill-of-materials references. Remote attestation may be used where the deployment can establish a meaningful hardware/software root of trust.

\subsection{Monotone Safety Updates}
Let \(P_2\preceq_{safe}P_1\) mean that every effect prohibited by \(P_1\) remains prohibited by \(P_2\), every safety invariant of \(P_1\) is preserved, and no new capability or exception broadens the admitted effect surface. Such updates are monotone restrictive and may use the ordinary safety-update quorum. Any change that can enlarge the admitted set, relax an invariant, or alter the constitution is classified as a new safety constitution and requires elevated independent approval, new proof artifacts, regression testing, and a new bundle hash.

\subsection{Incident Response}
Incident response is part of the assurance contract:
\[
\text{Detect}\rightarrow\text{Freeze}\rightarrow\text{Revoke}\rightarrow\text{Quarantine}\rightarrow\text{Preserve Evidence}\rightarrow\text{Revalidate}\rightarrow\text{Restore}.
\]
Freeze and revoke are monotone restrictive. Quarantine removes external effects. Evidence preservation is independent of the agent where the profile requires it. Restoration requires a new admitted \SafetyBundle{} or a new validated state, not merely a model restart.

\subsection{Low-Latency Integration}
The target \(\latTarget\le 1\,\mathrm{ms}\) applies only to the mechanical/schema path measured on the deployment's reference hardware and workload. It is not a promise that semantic extraction, human review, or full policy compilation can occur in one millisecond. High-risk decisions should be architected asynchronously where necessary, with pre-approved plans and state-bound capabilities so that latency does not force semantic shortcuts.

\subsection{Integration and Plan-Scoped Capabilities}
Before untrusted external data is read, the trusted principal's request is converted into an approved plan envelope for the expected effect set, resource scope, budget, and lifetime. Capabilities are then issued against that plan rather than against arbitrary future intents generated after the agent has consumed untrusted content. This reduces prompt-injection-induced privilege expansion and provides an explicit object against which delegation can be checked.

\subsection{Causal Continuation and Harm-Coverage Invariant}
Every harm class named in the Safety Profile must have at least one protection node on a mechanically covered channel if the deployment claims prevention for that class. If a harm class is represented only by semantic monitoring, it is assigned Tier~2 or monitored-only status. This rule prevents a paper-level harm taxonomy from being mistaken for an implementation-level causal guarantee.

\subsection{What the Runtime Architecture Does Not Solve}
No reference monitor described here prevents all possible human persuasion, covert channels, malicious physical access, or deliberate destruction of the monitor by an operator who controls the entire root of trust. No sensor contract proves an objective semantic category when that category is inherently contested. No cryptographic artifact by itself establishes legal liability. The strength of the framework is therefore proportional to the precision and independence of the deployment profile.

\section{Assurance, Coverage, and Measurement}\label{sec:guarantees}
\subsection{Four Assurance Classes}
The framework uses four primary claim classes.
\begin{longtable}{@{}L{2.3cm} L{5.0cm} L{7.1cm}@{}}
\toprule
\textbf{Class} & \textbf{Meaning} & \textbf{Evidence required}\\
\midrule
\endfirsthead
\toprule
\textbf{Class} & \textbf{Meaning} & \textbf{Evidence required}\\
\midrule
\endhead
\bottomrule
\endfoot
G-Normative-Commitment & Operator consistency with the selected constitution and its declared practice & Dependency graph, explicit constitution, claims registry, governance record. \\
G-Safety-Execution (\Gmech{}) & Conditional safety of mechanically covered protected effects & Verified kernel semantics, complete-mediation evidence, policy refinement proof, capability proof, state-revalidation evidence. \\
G-Semantic (\Gsem{}) & Empirical assurance for probabilistic semantic judgments & Pre-registered calibration, coverage/selective-risk analysis, adaptive red-team evaluation, drift monitoring. \\
G-Governance & Institutional controls governing updates, accountability, evidence retention, and emergency response & Independent approvals, transparency records, key custody, incident drills, audit reconstruction. \\
\end{longtable}

None of these classes subsumes another. A deployment can have strong G-Safety-Execution evidence while still having unresolved semantic or governance residuals. Conversely, a polished governance program cannot create a mechanical guarantee for an effect that is not actually mediated.

\subsection{Conditional Execution Guarantee}\label{sec:go}
Under C1--C8 and a valid \SafetyBundle{}, the protected-effect guarantee is:
\[
\forall s,a,s',\qquad \mathrm{Covered}(a)\land\Execute(s,a,s')\Rightarrow\Safe(s,a,s').
\]
This is a theorem schema. Its proof burden lies in the deployment artifacts that instantiate C1--C8. It is not an empirical claim that current implementations already satisfy the conditions.

\subsection{Categorical versus Probabilistic Semantic Guarantees}
The semantic layer is categorical only when the semantic predicate is represented by deterministic, independently verifiable data. For probabilistic classes, the deployment reports selective risk and coverage. Let \(p_{high}\) be the measured false-allow probability for Tier~1 semantic decisions under the preregistered evaluation regime. For a horizon of \(N\) mediated decisions, the conservative union-bound report is
\[
P(\text{at least one semantic false-allow})\le Np_{high}
\]
without assuming independent actions. If the measured events are dependent, tighter bounds may be reported only when the dependence model is separately justified. High-consequence Tier~0 effects are not authorized from this probability alone.

\subsection{Residual Risk Budget}
Each deployment declares a Residual Risk Budget containing, at minimum, residual channels, expected exposure, severity category, detection capability, response time, compensating controls, and owner. The budget is not an excuse to leave gaps unnamed. It is the accounting mechanism for channels that remain outside the categorical theorem. A deployment whose residual-risk budget is exceeded must automatically downgrade, freeze, or re-review the affected profile.

\newpage
\subsection{Coverage Matrix}\label{sec:gc}
The deployment publishes a matrix over effect classes and channels:
\begin{longtable}{@{}L{3.6cm} L{3.0cm} L{3.2cm} L{4.2cm}@{}}
\toprule
\textbf{Effect class} & \textbf{Protection tier} & \textbf{Primary control} & \textbf{Residual / evidence}\\
\midrule
\endfirsthead
\toprule
\textbf{Effect class} & \textbf{Protection tier} & \textbf{Primary control} & \textbf{Residual / evidence}\\
\midrule
\endhead
\bottomrule
\endfoot
File mutation & 0 & OS / gateway mediation + capability & Alternate process paths closed; audit witness. \\
Network egress & 0/1 & Egress gateway + IFC & Unmanaged interfaces listed or closed. \\
Credential use & 0 & Key service + sender-constrained capability & Replay and substitution tests. \\
Financial transfer & 0 & Typed authorization + bounded capability & Independent approval and result reconciliation. \\
Model/weight export & 0 & Storage / network mediation + provenance & Exfiltration red-team. \\
Inter-agent delegation & 1 & Attenuated capability lineage & Proxy and scope-escalation tests. \\
Communication to humans & 2 or monitored-only & Semantic monitor / typed template & Persuasion and covert-channel residuals declared. \\
Oversight control & 0 & Independent safety authority & Freeze / revoke path tested. \\
\end{longtable}

\subsection{Enterprise Audit Matrix}
\begin{longtable}{@{}L{3.2cm} L{4.0cm} L{4.0cm} L{2.8cm}@{}}
\toprule
\textbf{Control} & \textbf{Question} & \textbf{Evidence} & \textbf{Status}\\
\midrule
\endfirsthead
\toprule
\textbf{Control} & \textbf{Question} & \textbf{Evidence} & \textbf{Status}\\
\midrule
\endhead
\bottomrule
\endfoot
Complete mediation & Can any protected effect bypass the monitor? & Channel inventory, negative-space tests, platform controls & E/S. \\
Policy refinement & Does the deployed policy refine the safety constitution? & Proof object + checker output & S. \\
DSC soundness & Does the abstract state preserve safety? & Abstraction theorem / model checking & S/E. \\
Capability lineage & Can a child broaden authority? & Verified attenuation + tests & S. \\
TOCTOU & Can state change invalidate an authorization? & Atomic revalidation tests & S/E. \\
Sensor assurance & What is \(\epscover\) under the declared threat regime? & Held-out + adaptive evaluation & E. \\
Evidence integrity & Can events be suppressed or altered? & External witness / transparency log & S/E. \\
Update safety & Can an operator silently weaken the floor? & Quorum, monotonicity proof, bundle log & G. \\
\end{longtable}

\newpage
\subsection{Release Gates}
A Profile-1 release is not a single pass/fail benchmark. It is a sequence of ten gates.
\begin{enumerate}
\item \textbf{Gate 1 --- Formal Core:} machine-readable statements for L1, NRD, CR/CAR and the dependency graph.
\item \textbf{Gate 2 --- Safety Ontology:} closed effect, action, state, authorization, evidence, budget, and exception schemas.
\item \textbf{Gate 3 --- Safety Profile:} explicit protected effects, invariants, forbidden transitions, risk tiers, and fallback rules.
\item \textbf{Gate 4 --- Kernel Assurance:} verified policy semantics, refinement checker, capability verifier, and fail-safe matrix.
\item \textbf{Gate 5 --- Mediation Closure:} channel inventory plus negative-space and escape-path testing for the declared profile.
\item \textbf{Gate 6 --- Sensor Assurance:} held-out, adversarial, adaptive, and drift evaluation with published \(\epscover\) and selective risk.
\item \textbf{Gate 7 --- Supply Chain:} signed \SafetyBundle{}, model/compiler/kernel hashes, dependency manifest, and reproducible build evidence.
\item \textbf{Gate 8 --- Governance:} independent approval, key custody, update rules, residual-risk budget, incident drills, and public nutrition label.
\item \textbf{Gate 9 --- Independent Challenge:} adversarial review by a party that did not author the deployment controls, including long-horizon and proxy attacks.
\item \textbf{Gate 10 --- Release or Downgrade:} publish the exact claim scope. Missing evidence automatically downgrades the claim class rather than being hidden in prose.
\end{enumerate}

\subsection{Evaluation Strata}
Evaluation is organized into four strata: public benchmarks; generated adversarial policy tests; deployment-specific red-team scenarios; and held-out adaptive evaluation whose attack strategies are not available to the developers during the final tuning stage. Relevant public benchmarks include AgentDojo and AgentHarm; dynamic prompt-injection evaluation is also represented by AgentDyn \cite{agentdojo,agentharm,agentdyn}. These benchmarks are evidence of challenge coverage, not proofs of universal safety.

\subsection{Severity-Stratified and Horizon-Level Metrics}
Aggregate unauthorized-action rate is insufficient. The core metrics are:
\[
R_{UA,high},\quad R_{UA,1},\quad E[Loss]/10^6\ \text{decisions},\quad M_{effect},\quad P(\text{no catastrophic failure over }N),
\]
plus false-block rate, legitimate task completion, audit-reconstruction time, evidence completeness, and p99 mechanical overhead. Results are reported by severity, effect class, principal, attack family, and policy version. Maximum observed effect magnitude is retained even when the event rate is small.

\subsection{Dependence-Aware Statistical Discipline}
The default safety calculation does not assume i.i.d. actions. Sequence-level tests use blocking by task, principal, policy version, and attack family. Adaptive attackers are treated as part of the evaluation regime. Where multiple evidence channels contribute to a bound, dependence is reported rather than ignored. A union bound is the safe default; tighter composition is treated as an additional claim requiring a justified dependence model.

\subsection{Pre-Registered Ablation}
The ablation has at least four arms: baseline without the governor; containment-only; full governor; and component-level variants that selectively remove policy refinement, semantic fail-safe extraction, capability lineage, trajectory monitoring, evidence witnessing, or constitutional content. Co-primary outcomes include unauthorized protected execution, legitimate task completion, p99 mechanical overhead, audit-reconstruction time, compliance-artifact completeness, and independent evidence quality. The experiment is not considered complete until its analysis plan is frozen before final evaluation.

\subsection{Formal Closure and Reproducibility}
The machine-checking program must publish: formal statements of L1, NRD, and CR/CAR; the typed \Ground{} calculus; the \ADM{}/NRC enumeration checker; formal safety-profile schemas; policy-refinement checker code; environment manifests; seeds or exhaustive parameters; proof outputs; and the exact compiler and kernel versions used for the release. The paper does not claim these artifacts exist merely because it specifies them.

\subsection{Evidence and Implementation Status}\label{sec:status}
This paper is a complete conceptual and engineering specification. The philosophical dependency structure is explicit; the main security theorems are stated conditionally; the deployment profile and \SafetyBundle{} schema are specified; and the evaluation protocol is defined. What remains empirical or implementation-dependent is deliberately labeled: machine-checked proof of the philosophical derivations; implementation of the reference monitor; adversarial sensor assurance; complete-mediation evidence for any concrete platform; cryptographic key-custody testing; policy compiler fidelity; incident-response drills; and independent red-team validation.

\subsection{The Honesty Ledger}\label{sec:limits}
The following statements are prohibited as summaries of the framework:
\begin{enumerate}
\item ``Morality follows from Identity.''
\item ``Every rational AI agent is therefore morally bound.''
\item ``The model accepts the constitution, so the execution path is safe.''
\item ``A calibrated semantic sensor establishes categorical safety.''
\item ``A signed log proves legal non-repudiation or liability.''
\item ``A benchmark with zero observed failures proves zero real-world risk.''
\item ``Complete mediation is universal rather than deployment-specific.''
\end{enumerate}
The corresponding permitted claims are narrower: the philosophy specifies and audits a selected constitution; the \SafetyBundle{} makes implementation commitments explicit; and the reference monitor can provide a conditional execution theorem on the effect surface for which complete mediation and the other stated assumptions are actually demonstrated.

\subsection{Residual Region}
The residual region includes unmediated physical effects; hidden or covert channels; persuasive communication whose safety criterion is not objectively labelable; compromised root-of-trust infrastructure; operator-controlled disabling outside the governed profile; semantic categories whose truth cannot be independently verified; and future attack strategies outside the evaluation regime. These are not failures of terminology. They are the boundary of the theorem.

\Needspace{12\baselineskip}
\begin{keyresult}
\textbf{Core claim hierarchy.}
\[
\begin{gathered}
\mathrm{D:}\;L1,\;NRD,\;CR/CAR\\
\mathrm{A:}\;B4\text{--}B8,\;\text{Safety Profile}\\
\mathrm{S:}\;\text{kernel, refinement, and mediation invariants}\\
\mathrm{E:}\;\text{sensor and empirical assurance}\\
\mathrm{R:}\;\text{residual channels and pending evidence}
\end{gathered}
\]
A publication-ready system should preserve these labels in machine-readable release metadata.
\end{keyresult}

\medskip

\textbf{Closing.} The framework is strongest when its two ambitions are kept distinct and then connected through explicit interfaces. The practice-based ethics gives a transparent account of reason-giving and the declared constitution gives substantive content. The safety profile turns that content into a typed machine contract. The reference monitor turns the contract into an execution boundary. The assurance and governance layers state when the boundary is actually demonstrated and what remains outside it. No layer is used as a substitute for another. This is what permits the framework to make a serious enterprise claim without requiring a fictional guarantee: when the protected effect surface is closed, the proof obligations hold, and the deployment artifacts are independently checkable, unauthorized protected execution is excluded by the enforcement architecture rather than by the goodwill of the model.

\clearpage
\appendix
\section{The Compressed Ethical Chain}\label{app:chain}
For the rigorous reader, the ethical chain is presented in compressed form. The chain uses only the premises named in this paper.
\[
\begin{array}{rcl}
\mathrm{P1} & : & A = A \text{, non-contradiction} \\
& \downarrow & \text{(does not entail ethics)} \\
\mathrm{B1} & : & \text{\Atwo{} (Addressed Justification)} + \text{ARC} \\
+\ \mathrm{B2} & : & \text{RR (Responsiveness / Discharge)} \\
& \rightarrow & \text{Accountable reason-giving (Level 2)} \\
& \rightarrow & \text{Universalization of stated grounds (L1, \ca{})} \\
& \rightarrow & \text{T1 (Ground-Indexed Consistency), CE, CSL} \\
\mathrm{B3} & : & \text{AR (Answerability for Normative Difference)} \\
& \rightarrow & \text{NRD (No Unjustified Normative Difference, \ca{})} \\
& \rightarrow & \text{CR / CAR (Constitutional Reflexivity / Answerability, meta-level)} \\
\mathrm{B4} & : & \text{Agency-Completeness} \\
& \rightarrow & \text{F2 (Agency-Reciprocity)} \\
\mathrm{B5} & : & \text{Defensive Interpretation} \\
& \rightarrow & \text{F3 (Defensive Content)} \\
\mathrm{B6} & : & \text{A3 and Basic Goods (a--d)} \\
\mathrm{B7} & : & \text{PRAC (idealization)} \\
\mathrm{B8} & : & \text{Forward-Looking Grounds} \\
& \rightarrow & \text{T2-D, T2-G, T3, T4, T5, T6, T7a, T8, T11a, T16}
\end{array}
\]
\textbf{Status of the chain.} L1, NRD, and CR/CAR are the framework's non-trivial derivations (\ca{}; CR/CAR at the meta-level). All other results are \cb{} restatements of the basic constitutive standards, or contingent extensions. The CR/CAR derivation is a proof sketch; like the rest of the chain, it has not been machine-checked.

\textbf{Enforcement (architectural).} NO VALID CAPABILITY $\Rightarrow$ NO PROTECTED EFFECT, under C1--C8 for mechanically covered channels. The speech channel is covered by \Gsem{} only to the extent that sensor adequacy (SA) is established for the relevant semantic class.

\textbf{Containment (architectural).} Plane separation; mechanical enforcement; assume breach; budgets on principals; least privilege; and the boundary between the governed agent and its authority domain treated as a boundary rather than as a negotiable norm.

\section{Dependency-Hygiene Criteria}\label{app:routeo}
This specification is closed by default: it admits no premise, derivation, or interpretation from outside the dependency graph of Sections~\ref{sec:thin}--\ref{sec:rab}. Any future extension that proposes to derive \Atwo{} + ARC, or B3--B8, from an external ontological or metaphysical framework is admitted only if it meets all of the following criteria.
\begin{enumerate}
\item \textbf{H-1, explicit chain.} The derivation is stated step by step from its own premises alone, with no appeal to metaphor, analogy, or relabeling.
\item \textbf{H-2, named normative bridge.} Every normative term in the conclusion is either derived from non-normative premises or the bridge is declared as a premise with a status and a stated cost of rejection.
\item \textbf{H-3, compossibility.} The derivation is consistent with the observation that violations of the derived rule occur.
\item \textbf{H-4, inherited status.} The external premises used carry their own declared status, and the derived result inherits the weakest of them.
\item \textbf{H-5, independent checking.} The derivation survives machine-checking and independent adversarial review under the same protocol as Gate~1 and Gate~9 of Section~\ref{sec:guarantees}.
\end{enumerate}
No ontological import, exception, or semantic relabeling is admitted without passing these criteria.

\section{Register and Dependency Audit}\label{app:audit}
\textbf{Method.} The framework is transcribed into a register of more than seventy items, each with a class, a status, and its declared dependencies. A checking script performs the following operations:
\begin{enumerate}
\item Every cited dependency is a declared item.
\item The dependency graph contains no cycles.
\item The transitive set of non-derived premises under each derived result is computed.
\item Any item labeled D that rests on a C or U premise is flagged.
\item The results that fall when each premise is removed are listed.
\item No register item depends on an interpretive gloss rather than on a declared premise.
\end{enumerate}
\textbf{Results.}
\begin{itemize}
\item \textbf{Undeclared dependencies:} none in the register.
\item \textbf{Cycles:} none.
\item \textbf{Interpretive dependencies:} none.
\item \textbf{Single-premise dependency ablation over the core results:} removing B1 or B2 removes all core results; removing B3 removes NRD, CR/CAR and the B3-dependent extensions; removing B4 removes F2 and the B4-dependent general-agency results; removing B5 removes F3, T4, T5 and the B5-dependent results; removing B6 removes the B6-dependent protection, necessity, and conflict-ordering results; removing B7 removes the results that explicitly use PRAC (including T2-D, T3, T4, and T5); removing B8 removes the B8-dependent non-retribution result. This is a dependency report, not an independent validity proof.
\end{itemize}
\textbf{Limits of the audit.} The register is a transcription. The audit checks the structure of the declared derivations. It does not check that any derivation is valid, that the premises are consistent, or that they are independent. Those are tasks for the machine-checking program of Section~\ref{sec:guarantees} (Gate~1).

\section{Literature Positioning}\label{app:literature}
The framework's structure is adjacent to several established traditions: Gewirth's argument from the necessary conditions of agency \cite{gewirth}; Darwall's second-person standpoint and the account of accountability it supports \cite{darwall}; Scanlon's contractualist appeal to principles no one could reasonably reject \cite{scanlon}; Habermas's discourse ethics \cite{habermas}; Apel's transcendental pragmatics \cite{apel}; and Brandom's normative pragmatics of discursive commitment \cite{brandom}. The framework's claimed contribution is not to improve on any of these arguments but to \emph{package} a constitutive-presupposition strategy in a specific form: explicit premises, published dependency graph, honest costs of rejection, ablatable rules, machine-checkable structure, version discipline, and honest limits.

\section{AI Safety Literature Positioning}\label{app:ai-safety}
The framework's runtime architecture and governance mechanisms are adjacent to a body of AI safety work.
\subsection{Guaranteed-Safe AI and Runtime Verification}
\Gmech{} is a specification for a class of systems whose mechanically covered actions are gated by a credential-holding kernel \cite{seshia,bartocci-falcone}. Recent agent-security work shows that model-based monitors can themselves be bypassed through agent-as-proxy constructions, reinforcing the paper's decision to keep the final capability boundary external to the monitor's natural-language judgment \cite{agentproxy}.

\subsection{Capability Security and the Confused Deputy}
The capability-based access model is grounded in capability-security ideas, including the confused-deputy problem \cite{hardy}. The specified capability artifact adds cryptographic binding and expiry so that provenance, replay resistance, and principal attribution become testable properties rather than assumptions.

\subsection{Constitutional Controller and FormalJudge}
The Constitutional Controller (CoCo) treats a constitution as a first-class artifact and combines probabilistic and symbolic control \cite{coco}. The framework's runtime architecture shares this general neuro-symbolic direction but adds constitutional reflexivity, the split kernel, uncertainty-bounded state extraction, signed capabilities, and explicit assurance conditions (C1--C8 and SA). FormalJudge decomposes natural-language requirements into formal constraints and checks them with formal methods including SMT solving \cite{formaljudge}. The framework's mechanically covered channels are analogous to this formalized oversight boundary; the speech-channel guarantee (\Gsem{}) is not analogous, because it remains sensor-dependent.

\subsection{Agentic Security Benchmarks and Assurance Standards}
The empirical plan is intentionally multi-layered. AgentHarm evaluates harmful agent behavior; AgentDojo evaluates indirect prompt injection; and AgentDyn stresses open-ended dynamic evaluation beyond fixed static cases \cite{agentharm,agentdojo,agentdyn}. The deployment threat matrix additionally draws on the OWASP Top 10 for Agentic Applications \cite{owasp-agentic}. The governance interface maps to Article~14 of the EU AI Act (human oversight) \cite{euaiact}, ISO/IEC 42001:2023 (AI management systems) \cite{iso42001}, and NIST AI RMF 1.0 plus its Generative AI Profile \cite{nistairmf,nistgenai}. The compiler treats these sources as inputs to a traceable obligation-to-control process, not as documents that can be mechanically converted into law.

\section{Empirical Evaluation Protocol}\label{app:evaluation}
\subsection{Primary Safety Outcomes}
The primary outcomes are unauthorized protected execution stratified by severity, maximum observed effect magnitude, cumulative loss per million decisions, horizon-level catastrophic-failure probability, and channel-bypass rate. Secondary outcomes include false-block rate, legitimate task completion, p99 mechanical overhead, audit-reconstruction time, evidence completeness, and recovery time after freeze/revoke.

\subsection{Evaluation Strata}
The benchmark suite is partitioned into: public benchmarks; generated adversarial policies; deployment-specific red-team scenarios; and held-out adaptive evaluations. Attack classes include prompt injection, tool misuse, credential replay, principal substitution, proxy delegation, TOCTOU, policy downgrade, model update drift, multi-step decomposition, cumulative budget attacks, hidden subprocesses, alternate network paths, evidence suppression, colluding reviewers, emergency-control abuse, and supply-chain substitution.

\subsection{Sensor Assurance Protocol}
For every semantic predicate used in a categorical or probabilistic decision, preregister the ontology, label definition, abstention rule, \(\Gamma\)-construction method, coverage target, adversary class, perturbation budget, shift regime, and drift trigger. Report coverage and selective risk by class, tier, principal, horizon, and attack family. A sensor is not considered assured merely because average calibration is good.

\subsection{Horizon-Level Analysis}
For a sequence of \(N\) decisions, compute the probability of no catastrophic failure over the complete horizon, not only per-decision error. Per-decision rates are retained for diagnosis. The paper's default analytical bound uses a union bound; independence assumptions require separate evidence.

\subsection{Ablation Matrix}
The evaluation must compare at least the following:
\begin{longtable}{@{}L{4.4cm} L{9.2cm}@{}}
\toprule
\textbf{Arm} & \textbf{Purpose}\\
\midrule
\endfirsthead
\toprule
\textbf{Arm} & \textbf{Purpose}\\
\midrule
\endhead
\bottomrule
\endfoot
No governor & Measures unconstrained agent behavior. \\
Containment only & Tests what credential and channel controls achieve without semantic/constitutional layers. \\
Full governor & Tests the complete stack. \\
No policy refinement & Tests the contribution of state-transition admission. \\
No capability lineage & Tests proxy and delegation exposure. \\
No trajectory monitor & Tests decomposition and cumulative-harm exposure. \\
No independent evidence witness & Tests evidence suppression resilience. \\
Constitutional variants & Tests whether different constitutions change security or utility only through the declared policy content. \\
\end{longtable}

\subsection{Statistical and Reproducibility Requirements}
All thresholds, random seeds, hardware, software versions, attack budgets, and exclusion criteria must be fixed before final evaluation. Report confidence intervals or other uncertainty quantification appropriate to the estimator. Where the attacker adapts to the system, the evaluation protocol must model that adaptation rather than assume i.i.d. trials. Full traces and machine-readable metrics should be retained for independent reconstruction.

\section{Deployment Roadmap}\label{app:roadmap}
The implementation path is organized around assurance closure rather than feature accumulation. Its ten steps correspond one-to-one to the release gates of Section~\ref{sec:guarantees}.
\begin{enumerate}
\item \textbf{Formal Core.} Machine-check the dependency graph, typed \Ground{} calculus, L1/NRD/CR/CAR statements, and the \ADM{}/NRC artifact.
\item \textbf{Safety Ontology.} Freeze schemas for principals, resources, actions, effects, states, histories, authorities, budgets, evidence, and exceptions.
\item \textbf{Safety Profile.} Instantiate the protected effect set, invariants, forbidden transitions, tiers, aggregation rules, fallback rules, and residual-risk budget for one concrete deployment profile.
\item \textbf{Kernel Assurance.} Implement the minimal TCB, proof checker, capability verifier, revocation path, and fail-safe matrix; formally test conformance.
\item \textbf{Mediation Closure.} Enumerate and test alternate channels, including subprocesses, plugins, schedulers, browser paths, inter-agent links, and external service proxies.
\item \textbf{Sensor Assurance.} Establish \(\epscover\), selective risk, disagreement handling, adaptive coverage, and drift triggers for every semantic predicate that remains in the protected path.
\item \textbf{Supply Chain.} Bind the \SafetyBundle{} to signed builds, dependency manifests, and model/compiler/kernel hashes, with transparency logging.
\item \textbf{Governance.} Establish threshold approval, key custody, update rules, residual-risk budget, incident response and drills, and the public nutrition label. Policy fidelity tooling (restricted DSL, dual translation or independent adjudication, mutation testing, example-set generation) is completed within the Kernel Assurance and Governance steps.
\item \textbf{Independent Challenge.} Conduct deployment-specific red-teaming, long-horizon evaluation, and independent proof/evidence review.
\item \textbf{Release or Downgrade.} Publish exactly the claim class supported by the collected evidence. Any unmet gate downgrades the deployment's claim scope rather than being hidden by a general ``safe'' label.
\end{enumerate}

\section{Glossary of Named Items}\label{sec:glossary}
\begin{longtable}{@{}L{3.2cm} L{11.8cm}@{}}
\toprule
\textbf{Term} & \textbf{Definition} \\
\midrule
\endfirsthead
\toprule
\textbf{Term} & \textbf{Definition} \\
\midrule
\endhead
\bottomrule
\endfoot
P1 (Identity) & Identity premise used as the logical background; it does not by itself entail substantive ethics. \\
\Atwo{} (Addressed Justification) & A claim on another offers a justificatory reason to the addressee. \\
ARC & Addressee-Reason Condition: the justificatory ground must be addressable in the sense specified by the practice. \\
RR & Responsiveness: challenged justification must be defended, qualified, or withdrawn. \\
RR-G & No-Retrofitting Rule for typed grounds. \\
AR & Answerability for normative difference. \\
CR / CAR & Constitutional Reflexivity / Constitutional Answerability. \\
L1-Public & Invariant public-difference component of L1. \\
L1-Relevance & Relevance component of L1, supported by the stronger practice premises rather than \ADM{} alone. \\
NRD & No Unjustified Normative Difference. \\
NRC & Normative Role Closure: role alone may not alter the normative force of a ground. \\
FSP & Floor--Superstructure Principle: agent-relative duties may add to but never subtract from the floor owed to all (used by T8). \\
\Ground{} & Typed normative ground object with predicate, domain, evidence, scope, effect, and provenance. \\
DSC & Decision State Contract governing the machine-readable state supplied to the safety kernel. \\
Safety Profile & Versioned machine-readable specification of protected effects, invariants, transitions, tiers, budgets, and fallback rules. \\
\SafetyBundle{} & Cryptographically bound artifact linking constitution, profile, policy, proof, kernel, sensor, model, compiler, deployment profile, and dependencies. \\
Mediation Invariant & \(\neg\mathrm{ValidCapability}(e)\Rightarrow\neg\mathrm{Execute}(e)\) for covered effects. \\
\PRef{} & Typed policy-refinement relation between deployment policy and safety constitution. \\
\Safe{} & State-transition predicate defining admitted safety for a protected effect. \\
\Gmech{} & Mechanical execution assurance for covered effects, conditional on C1--C8. \\
\Gsem{} & Empirical assurance for probabilistic semantic judgments, conditional on SA. \\
G-Governance & Institutional assurance for updates, accountability, evidence, and incident response. \\
GC & Complete mediation / channel completeness (C1). \\
TCB & Trusted-core integrity (C2). \\
KVC & Kernel verification and conformance (C3). \\
PAV & Policy-artifact verification (C4). \\
PA & Principal attribution and capability non-transferability (C5). \\
SA & Sensor adequacy: the declared coverage and error bound of the semantic sensor holds in the evaluated regime. \\
\end{longtable}

\section*{Acknowledgments}

This research was conducted entirely independently, without institutional affiliation or external support. The author used AI language models as computational co-auditors and structural stress-testers, including for symbolic calculation, structural review, and document preparation. The author retains full responsibility for the concepts, arguments, citations, and final claims.

\newpage
\phantomsection
\addcontentsline{toc}{section}{References}
\begin{thebibliography}{99}
\CanonBibliographySetup

\bibitem{gewirth} Gewirth, A. \textit{Reason and Morality}. University of Chicago Press, 1978.
\bibitem{darwall} Darwall, S. \textit{The Second-Person Standpoint: Morality, Respect, and Accountability}. Harvard University Press, 2006.
\bibitem{scanlon} Scanlon, T. M. \textit{What We Owe to Each Other}. Belknap Press of Harvard University Press, 1998.
\bibitem{habermas} Habermas, J. \textit{Moral Consciousness and Communicative Action}. MIT Press, 1990.
\bibitem{apel} Apel, K.-O. \textit{Towards a Transformation of Philosophy}. Routledge, 1980.
\bibitem{brandom} Brandom, R. B. \textit{Making It Explicit: Reasoning, Representing, and Discursive Commitment}. Harvard University Press, 1994.
\bibitem{christiano} Christiano, P. F., Leike, J., Brown, T. B., Martic, M., Legg, S., and Amodei, D. ``Deep Reinforcement Learning from Human Preferences.'' \textit{Advances in Neural Information Processing Systems}, 30:4299--4307, 2017.
\bibitem{bai-constitutional} Bai, Y., Kadavath, S., Kundu, S., Askell, A., Kernion, J., Jones, A., Chen, A., Goldie, A., Mirhoseini, A., McKinnon, C., Chen, C., Olsson, C., Olah, C., Hernandez, D., Drain, D., Ganguli, D., Li, D., Tran-Johnson, E., Perez, E., Kerr, J., Mueller, J., Ladish, J., Landis, J., Ndousse, K., Nguyen, K., Askell, K., Amodei, D., Clark, J., McCandlish, S., et al. ``Constitutional AI: Harmlessness from AI Feedback.'' arXiv:2212.08073, 2022.
\bibitem{scalable-oversight} Bowman, S. R., Hyun, J., Perez, E., Chen, E., Pettit, C., Heiner, S., et al. ``Measuring Progress on Scalable Oversight for Large Language Models.'' arXiv:2211.03540, 2022.
\bibitem{coco} Kohaut, S., Divo, F., Hamid, N., Flade, B., Eggert, J., Dhami, D. S., and Kersting, K. ``The Constitutional Controller: Doubt-Calibrated Steering of Compliant Agents.'' arXiv:2507.15478, 2025.
\bibitem{formaljudge} Zhou, J., Sheng, Y., Lou, H., Yang, Y., and Fu, J. ``FormalJudge: A Neuro-Symbolic Paradigm for Agentic Oversight.'' In \textit{Proceedings of the 43rd International Conference on Machine Learning}, PMLR 306:165723--165752, 2026.
\bibitem{seshia} Seshia, S. A., Sadigh, D., and Sastry, S. S. ``Toward Verified Artificial Intelligence.'' \textit{Communications of the ACM}, 65(7):46--55, 2022. doi:10.1145/3503914. Earlier technical report: arXiv:1606.08514.
\bibitem{bartocci-falcone} Bartocci, E. and Falcone, Y., eds. \textit{Lectures on Runtime Verification: Introductory and Advanced Topics}. Lecture Notes in Computer Science, vol. 10457. Springer, 2018.
\bibitem{hardy} Hardy, N. ``The Confused Deputy (or Why Capabilities Might Have Been Invented).'' \textit{ACM SIGOPS Operating Systems Review}, 22(4):36--38, 1988. doi:10.1145/54289.871709.
\bibitem{agentharm}
Andriushchenko, M., Souly, A., Dziemian, M., Duenas, D., Lin, M., Wang, J., Hendrycks, D., Zou, A., Kolter, J. Z., Fredrikson, M., Gal, Y., and Davies, X. ``AgentHarm: A Benchmark for Measuring Harmfulness of LLM Agents.'' \textit{International Conference on Learning Representations (ICLR)}, 2025.
\bibitem{agentdojo}
Debenedetti, E., Zhang, J., Balunovic, M., Beurer-Kellner, L., Fischer, M., and Tramer, F. ``AgentDojo: A Dynamic Environment to Evaluate Prompt Injection Attacks and Defenses for LLM Agents.'' \textit{Advances in Neural Information Processing Systems}, 37, Datasets and Benchmarks Track, 2024. doi:10.52202/079017-2636.
\bibitem{agentdyn} Li, H., Wen, R., Shi, S., Zhang, N., Vorobeychik, Y., and Xiao, C. ``AgentDyn: Are Your Agent Security Defenses Deployable in Real-World Dynamic Environments?'' arXiv:2602.03117, 2026. doi:10.48550/arXiv.2602.03117.
\bibitem{agentproxy} Isbarov, J. and Kantarcioglu, M. ``Bypassing AI Control Protocols via Agent-as-a-Proxy Attacks.'' arXiv:2602.05066, 2026. doi:10.48550/arXiv.2602.05066.
\bibitem{rfc7515} Jones, M., Bradley, J., and Sakimura, N. ``JSON Web Signature (JWS).'' RFC 7515, IETF, 2015. doi:10.17487/RFC7515.
\bibitem{rfc7519} Jones, M., Bradley, J., and Sakimura, N. ``JSON Web Token (JWT).'' RFC 7519, IETF, 2015. doi:10.17487/RFC7519.
\bibitem{rfc8705} Campbell, B., Bradley, J., Sakimura, N., and Lodderstedt, T. ``OAuth 2.0 Mutual-TLS Client Authentication and Certificate-Bound Access Tokens.'' RFC 8705, IETF, 2020. doi:10.17487/RFC8705.
\bibitem{rfc9449} Fett, D., Campbell, B., Bradley, J., Lodderstedt, T., Jones, M., and Waite, D. ``OAuth 2.0 Demonstrating Proof of Possession (DPoP).'' RFC 9449, IETF, 2023. doi:10.17487/RFC9449.
\bibitem{rfc9162} Laurie, B., Messeri, E., and Stradling, R. ``Certificate Transparency Version 2.0.'' RFC 9162, IETF, 2021. doi:10.17487/RFC9162
\bibitem{owasp-agentic} OWASP GenAI Security Project. ``OWASP Top 10 for Agentic Applications for 2026.'' December 9, 2025.
\bibitem{nistairmf} Tabassi, E. \textit{Artificial Intelligence Risk Management Framework (AI RMF 1.0)}. NIST AI 100-1, National Institute of Standards and Technology, 2023. doi:10.6028/NIST.AI.100-1.
\bibitem{euaiact}
European Union. ``Regulation (EU) 2024/1689 of the European Parliament and of the Council of 13 June 2024 laying down harmonised rules on artificial intelligence (Artificial Intelligence Act).'' \textit{Official Journal of the European Union}, consolidated version of 27 July 2026. Art. 14, Human Oversight.
\bibitem{iso42001} ISO/IEC. \textit{ISO/IEC 42001:2023 Information technology --- Artificial intelligence --- Management system}. International Organization for Standardization, 2023.
\bibitem{nistgenai} Autio, C., Schwartz, R., Dunietz, J., Jain, S., Stanley, M., Tabassi, E., Hall, P., and Roberts, K. \textit{Artificial Intelligence Risk Management Framework: Generative Artificial Intelligence Profile}. NIST AI 600-1, National Institute of Standards and Technology, 2024. doi:10.6028/NIST.AI.600-1.
\bibitem{nist-agent-identity}
Booth, H., Fisher, W., Galluzzo, R., and Roberts, J. \textit{Accelerating the Adoption of Software and Artificial Intelligence Agent Identity and Authorization}. NIST National Cybersecurity Center of Excellence, Concept Paper, Initial Public Draft, February 5, 2026.
\end{thebibliography}

\end{document}